% !TEX root = ../main.tex
\chapter{All spacetimes beyond Einstein (Obergurgl lectures)}\label{ch:c29}
\newcommand{\obg}[1]{{\footnotesize\textsf{[Obergurgl, #1]}}}
\newcommand{\Dfib}{D}
\begin{paracol}{2}

\begin{sources}
\begin{tabularx}{\linewidth}{@{}l l L@{}}
\toprule
\textbf{Source} & \textbf{Length} & \textbf{Content}\\
\midrule
\href{https://arxiv.org/abs/1111.4824}{arXiv:1111.4824} & 44 pp. & F.~P.~Schuller, ``Obergurgl Lectures on All Spacetimes Beyond Einstein'' (2011). Written notes of eight lectures given in Obergurgl (Austria), first week of September 2011, twelve hours, at the summer school of the Elitestudiengang Physik (Erlangen, Nuremberg, Regensburg). I--II manifolds and tensors (a recall); III hyperbolic polynomials and the geometries that can carry predictive matter; IV massless dispersion, the dual polynomial and the Gauss map; V the definition of a spacetime (bi-hyperbolic, time-oriented, energy-distinguishing); VI massive dispersion, the Legendre map and the particle action; VII superluminal particles and why they are not seen; VIII spacetime dynamics: deformation algebra and the linear PDE system for the gravitational Lagrangian.\\
\bottomrule
\end{tabularx}

\smallskip
This chapter is built from a \emph{text}, not from a video: there are no video timestamps. Only one displayed formula of the paper is numbered (the cone condition of Lecture VII, eq.~(1)); otherwise a reference reads \obg{Lecture, item}: the lecture (I--VIII, as in the paper) and the name of the definition, theorem, remark or example. The paper is the written version of the lectures; it is based on D.~R\"atzel, S.~Rivera and F.~P.~Schuller, ``Geometry of physical dispersion relations'', Phys.\ Rev.\ D, 2010, as cited in the paper [A] and on K.~Giesel, F.~P.~Schuller, C.~Witte and M.~Wohlfarth, ``Dynamics of tensorial spacetime geometries'' (in preparation at the time, [B]). The text below paraphrases the paper in the paper's order; formulas are reproduced as printed, and where the printed formula looks wrong this is said in a note. The figures of the paper (zero sets, cones, the 1+1 dimensional momentum diagrams of Lecture VII) are not copied; those that matter are described in words or redrawn.
\end{sources}

\begin{reading}[Where this chapter sits in the other notes]
The programme is that of lecture~26 (chapter~\ref{ch:c26}), where the same objects (principal polynomial, dual polynomial, hyperbolicity, Gauss and Legendre map, deformation algebra, the PDE system) are introduced in words; here they are written as definitions, theorems and proofs, with the formulas of the lecture slides. The later papers are chapters~\ref{ch:c30} (Schuller and Witte, arXiv:1402.6548) and~\ref{ch:c31} (the review). Background in this book: manifolds, charts, atlases chapters~\ref{ch:c01}, \ref{ch:c02}, \ref{ch:c04}; tangent and cotangent spaces chapter~\ref{ch:c05}; tensors chapter~\ref{ch:c03}; bundles and tensor fields chapter~\ref{ch:c06}; metric manifolds chapter~\ref{ch:c10}; spacetime and matter chapters~\ref{ch:c13}, \ref{ch:c14}; the Einstein action chapter~\ref{ch:c15}; hypersurface deformations and the constraints of canonical gravity chapters~\ref{ch:c18} and~\ref{ch:c19} (not written yet). Other Schuller sets: \emph{Geometric Anatomy of Theoretical Physics} (differentiable manifolds, tensor bundles) and the \emph{Theoretische Mechanik} notes for the Poisson brackets of phase space. The notation is that of the paper.
\end{reading}

\section{Introduction: which tensors can be a spacetime?}\label{sec:c29-intro}

The paper opens with the OPERA claim (2011) of neutrinos faster than light, which it takes as a prompt: spacetime need not carry a Lorentzian metric, so which other tensor field could it carry? \obg{Introduction} The answer it announces is complete. Call a tensor field admissible if matter equations on it can be solved forward from initial data and if every observer assigns the same sign to the energy of a particle. Then the matter dispersion relation must be given by a totally symmetric, contravariant tensor field $P$ of even rank, and $P$ must satisfy three algebraic requirements: hyperbolicity, time-orientability and the energy-distinguishing property. Nothing more is needed to transfer the whole kinematics of Lorentzian spacetimes to such a geometry. Causality survives even if some matter propagates faster than the slowest light; to see why, convex analysis, real algebraic geometry and modern PDE theory have to be used together. Since this class exhausts the possible spacetime structures on a smooth manifold, it is the only one that deserves study.

A Lorentzian metric is the simplest member. Using a finer geometry removes the basis of the Einstein equations as we know them, but not the principle behind them found by the Wheeler school: \emph{gravity is the rule that carries a spatial geometry from one initial-data hypersurface to a neighbouring one, in such a way that the stack of all these spatial geometries is an admissible spacetime geometry.} The principle is blind to the kind of spacetime geometry it is applied to. The second message is therefore: dynamics for a refined spacetime geometry need no inspired guess, they are the solution of a clean mathematical problem. Which of the (only countably many) classical spacetime theories is realised in nature is for experiment to decide.

\begin{coursenote}[How this chapter relates to lecture 26]
Lecture 26 (chapter~\ref{ch:c26}) gave the same story in the order: matter, principal polynomial, hyperbolic polynomials, spacetime, dynamics. The paper is the same story with the mathematics and the formulas, and with two extra topics: massless and massive point particles (Lectures IV and VI) and superluminal particles (Lecture VII).
\end{coursenote}

\section{Lecture I: Manifolds}\label{sec:c29-L1}

\obg{Lecture I} The foundation that is not questioned: spacetime is a smooth $d$-dimensional manifold $M$. The first lecture recalls the needed topology and manifold theory for the non-specialist. This is all in chapters~\ref{ch:c01}, \ref{ch:c02} and~\ref{ch:c04}; the definitions are only listed here, in the paper's order.

\begin{definition}[Topological space]\label{def:c29-top}
A set $M$ together with a family $\mathcal O$ of subsets, called open, such that (1) $\emptyset,M\in\mathcal O$; (2) $U,V\in\mathcal O\Rightarrow U\cap V\in\mathcal O$; (3) for any family $U_\alpha\in\mathcal O$, $\alpha\in A$, also $\bigcup_{\alpha}U_\alpha\in\mathcal O$.
\end{definition}
\begin{recall}[Topology]
Chapter~\ref{ch:c01}, \cref{def:c01-topology} and \cref{def:c01-open-closed}.
\end{recall}

Remarks of the paper: (1) the choice of topology on a set is far from unique; a set with 1, 2, 3, 4, 5, 6, 7 elements has 1, 4, 29, 355, 6942, 209527, 9535241 topologies, and the number grows quickly. (2) The coarsest topology is $\mathcal O=\{\emptyset,M\}$, the finest is the power set $\mathcal P(M)=\{U\mid U\subset M\}$; useful topologies lie in between (\cref{def:c01-extreme}). (3) Only for a finite set can one list the open sets; for an infinite set one defines them indirectly and then proves that they form a topology, as in the next definition.

\begin{definition}[Standard topology on $\R^d$]\label{def:c29-std}
$U\subset\R^d$ is open if for every $x\in U$ there is a positive real $\epsilon$ with $B_\epsilon(x)=\{y\in\R^d\mid\sum_{i=1}^d(x^i-y^i)^2<\epsilon\}\subset U$ (the paper writes the squared radius $\epsilon$ in the ball; this does not change the topology). Unless said otherwise $\R^d$ carries this topology.
\end{definition}

\begin{exercise}\label{ex:c29-rd}
Verify the three axioms for $(\R^d,\mathcal O_s)$ \obg{Lecture I, Exercise}. (Hint: for (2) take the smaller of two radii, for (3) the radius that works for one member of the union.)
\end{exercise}

\begin{definition}[Continuous map, homeomorphism]\label{def:c29-cont}
A map $\phi:M\to N$ of topological spaces is continuous if $\operatorname{preim}_\phi(V)=\{x\in M\mid\phi(x)\in V\}\in\mathcal O_M$ for all $V\in\mathcal O_N$. A bijection with $\phi$ and $\phi^{-1}$ continuous is a homeomorphism; spaces related by one are homeomorphic.
\end{definition}
\begin{recall}[Continuity]
\cref{def:c01-continuous}, \cref{def:c01-homeo}.
\end{recall}

Remarks: (1) every map is continuous if $N$ has the chaotic topology or $M$ the discrete one, another illustration that useful topologies lie between the extremes; (2) for $f:\R\to\R$ the definition reduces to the $\epsilon$-$\delta$ criterion.

\begin{definition}[Topological manifold, chart]\label{def:c29-manifold}
A topological space $(M,\mathcal O)$ is a $d$-dimensional topological manifold if every $p\in M$ has an open neighbourhood $U$ and a homeomorphism $\phi:U\to\phi(U)\subset\R^d$. The components $\phi^1,\dots,\phi^d$ are the coordinate functions, $(\phi^1(q),\dots,\phi^d(q))$ are the coordinates of $q\in U$ in the chart $(U,\phi)$.
\end{definition}
\begin{recall}[Manifolds]
\cref{def:c02-manifold}, \cref{def:c02-chart}.
\end{recall}

Remarks: (1) Charts are the sailor's way of working: a ship's route is a curve $\gamma:\R\to M$ on the real sea, and its chart representative $\gamma_\phi=\phi\circ\gamma:\operatorname{preim}_\gamma(U)\to\phi(U)$ is the line pencilled on the sea chart. (2) In general one chart does not suffice (maps of the Earth), one needs overlapping charts and a rule for passing between them.

\begin{definition}[$C^k$-compatibility, atlas, differentiable manifold]\label{def:c29-atlas}
Two charts $(U,\phi),(V,\psi)$ are $C^k$-compatible if either $U\cap V=\emptyset$, or $U\cap V\ne\emptyset$ and the transition map $\psi\circ\phi^{-1}:\phi(U\cap V)\to\psi(U\cap V)$ (a map between open subsets of $\R^d$) is invertible and, with its inverse, $k$ times continuously differentiable. A $C^k$-atlas is a collection $(U_\alpha,\phi_\alpha)_{\alpha\in A}$ of pairwise compatible charts whose domains cover $M$; it is maximal if it already contains every chart that is compatible with all its members. A topological manifold with a maximal $C^k$-atlas is a $C^k$-(differentiable) manifold.
\end{definition}
\begin{recall}[Atlas]
\cref{def:c02-transition}, \cref{def:c02-atlas}, \cref{def:c04-smooth}.
\end{recall}

Remarks of the paper: (1) A maximal $C^0$-atlas must lose charts to become a maximal $C^1$-atlas, and more to become $C^2$, and so on. In practice one gives a small $C^k$-atlas $\mathcal A_0$ and declares every chart compatible with all of $\mathcal A_0$ to belong to the maximal atlas $\mathcal A$. (2) A curve $\gamma$ is $m$ times continuously differentiable at $t_0$ if for some chart $(U,\phi)$ around $\gamma(t_0)$ the representative $\gamma_\phi$ is; the compatibility of the charts guarantees that this holds in any other chart $(V,\psi)$, because
\[
  \gamma_\psi=\psi\circ\gamma=\psi\circ\phi^{-1}\circ\phi\circ\gamma=(\psi\circ\phi^{-1})\circ\gamma_\phi,
\]
and the transition map is $k$ times continuously differentiable; so the definition does not depend on the chart. (3) Any two charts of a topological manifold are automatically $C^0$-compatible, so topological manifolds and $C^0$-manifolds are the same thing. Continuity can therefore be decided either on the topological manifold or on chart representatives, according to taste; differentiability, instead, can no longer be decided on the topological manifold: one must go down to the charts.

\begin{coursenote}[A slip in the paper]
The paper says $C^k$-differentiability ``for $k\geq0$'' cannot be decided at the topological level; since continuity ($k=0$) can, as the previous sentence says, it must mean $k\geq1$.
\end{coursenote}

\begin{exercise}\label{ex:c29-fun}
(a) Say what it means for $f:M\to\R$ to be $m$ times continuously differentiable on a $C^{k}$-manifold with $k\ge m$, and check that the answer does not depend on the chart. (b) Check that the $k$ times continuously differentiable functions form a real vector space $C^k(M)$ with pointwise operations \obg{Lecture I, Exercise}.
\end{exercise}

\begin{example}[The plane $\R^2$, Cartesian and polar chart]\label{ex:c29-plane}
$M=\R^2$ with the standard topology is a 2-dimensional topological manifold covered by the single chart $(U,\phi)$, $U=\R^2$, $\phi^1(x,y)=x$, $\phi^2(x,y)=y$. This one-chart atlas is trivially $C^\infty$-compatible with itself, so extending it to a maximal atlas gives a $C^\infty$-manifold $(\R^2,\mathcal O,\mathcal A_{\max})$. Another chart of $\mathcal A_{\max}$ is $V=\R^2\setminus\{(s,0)\mid s\in[0,\infty)\}$ with $\psi^1(x,y)=\sqrt{x^2+y^2}$, $\psi^2(x,y)=\arctan(y/x)$. The transition map is
\begin{gather*}
  \phi\circ\psi^{-1}:\R^+\times(0,2\pi)\to\R^2\setminus\{(0,s)\mid s\in[0,\infty)\},\\
  (\phi\circ\psi^{-1})(r,\varphi)=(r\cos\varphi,r\sin\varphi),
\end{gather*}
which is invertible and $C^\infty$; hence $(V,\psi)$ belongs to the maximal atlas \obg{Lecture I, Example}.
\end{example}
\begin{coursenote}[Polar chart as printed]
The paper's $\arctan(y/x)$ takes values in $(-\pi/2,\pi/2)$, while the transition map is stated on $\R^+\times(0,2\pi)$ and the removed set is the positive $x$-axis; the angle meant is the polar angle in $(0,2\pi)$ (and the image set listed has the roles of the axes swapped, it should be the plane minus the positive $x$-axis again). The point of the example is unaffected.
\end{coursenote}

\section{Lecture II: Tensors}\label{sec:c29-L2}

\obg{Lecture II} Without extra structure the only data one can put on a smooth manifold are tensor fields; the next lectures use them as prototypes of matter and of the fundamental geometry. Everything in this lecture is in chapters~\ref{ch:c03}, \ref{ch:c05} and~\ref{ch:c06}.

\begin{definition}[Tangent vector]\label{def:c29-tangent}
Let $M$ be a $C^\infty$-manifold and $\gamma$ a $C^1$-curve. The tangent vector $\dot\gamma_p$ to $\gamma$ at $p=\gamma(t_p)$ is the linear map $\dot\gamma_p:C^1(M)\to\R$, $\dot\gamma_pf=\dfrac{d(f\circ\gamma)}{dt}(t_p)$.
\end{definition}
\begin{recall}[Velocity of a curve]
\cref{def:c05-velocity}, \cref{def:c05-tangent}.
\end{recall}

Remarks of the paper. (1) The sum $(\dot\gamma_p\oplus\dot\delta_p)f=\dot\gamma_pf+\dot\delta_pf$ of two tangent vectors at $p$ is again a tangent vector, $\dot\sigma_p$ for the curve (built in a chart $(U,\phi)$ around $p$)
\[
  \sigma(t):=\phi^{-1}\bigl((\phi\circ\gamma)(t)+(\phi\circ\delta)(t)-(\phi\circ\gamma)(t_p)\bigr),
\]
since $\sigma(t_p)=p$ and by the chain rule, for any $C^1$-function $f$,
\[
  \dot\sigma_pf=\partial_a(f\circ\phi^{-1})(\phi(p))\Bigl(\frac{d\gamma_\phi{}^a}{dt}(t_p)+\frac{d\delta_\phi{}^a}{dt}(t_p)\Bigr)=\dot\gamma_pf+\dot\delta_pf.
\]
Likewise one defines the $S$-multiplication $\odot$ of a real number with a tangent vector. (2) The set $T_pM$ of tangent vectors at $p$ with $\oplus,\odot$ is a real vector space, the tangent space at $p$. (3) A chart $(U,\phi)$ induces a basis of $T_pM$: from
\[
  \dot\gamma_pf=\frac{d(\phi\circ\gamma)^a}{dt}(t_p)\,\partial_a(f\circ\phi^{-1})(\phi(p))
\]
and the definition $\frac{\partial}{\partial\phi^a}\big|_pf=\partial_a(f\circ\phi^{-1})(\phi(p))$ one gets $\dot\gamma_p=\frac{d\gamma_\phi{}^a}{dt}(t_p)\,\frac{\partial}{\partial\phi^a}\big|_p$: the components of a tangent vector in the chart-induced basis $\frac{\partial}{\partial\phi^1}\big|_p,\dots,\frac{\partial}{\partial\phi^d}\big|_p$ are the derivatives of the chart representative of the curve. (4) Under a change of chart $(U,\phi)\to(V,\psi)$,
\[
  \frac{\partial}{\partial\psi^a}\Big|_p=(C^\phi_\psi|_p)^m{}_a\frac{\partial}{\partial\phi^m}\Big|_p,\qquad (C^\phi_\psi|_p)^m{}_a=\partial_a(\phi\circ\psi^{-1})^m\big|_{\psi(p)},
\]
and the components change by $\frac{d\gamma_\psi{}^a}{dt}(t_p)=(C^\psi_\phi|_p)^a{}_m\frac{d\gamma_\phi{}^m}{dt}(t_p)$, where the two matrices are inverse to each other, $(C^\phi_\psi|_p)^m{}_a(C^\psi_\phi|_p)^a{}_n=\delta^m_n$.

\begin{definition}[Differential, cotangent space]\label{def:c29-diff}
The differential of a $C^1$-function $f$ at $p$ is the linear map $d_pf:T_pM\to\R$, $d_pf(X)=Xf$. The elements of the dual space $T^*_pM$ are the covectors.
\end{definition}
\begin{recall}[Cotangent space]
\cref{def:c05-cotangent}.
\end{recall}

Remarks. (1) The differentials $d_p\phi^1,\dots,d_p\phi^d$ of the chart components are elements of $T^*_pM$. (2) They act as $d_p\phi^a\bigl(\frac{\partial}{\partial\phi^b}\big|_p\bigr)=\partial_b(\phi^a\circ\phi^{-1})(\phi(p))=\delta^a_b$, so they are the dual basis of the chart-induced basis. (3) For $v=v_\phi{}^b\frac{\partial}{\partial\phi^b}\big|_p$ and $\omega=\omega^\phi{}_a\,d_p\phi^a$ one finds $\omega(v)=\omega^\phi{}_av_\phi{}^a$. (4) Under a change of chart, $d_p\psi^a=(C^\psi_\phi|_p)^a{}_b\,d_p\phi^b$.

\begin{definition}[Tensor of valence $(r,s)$ at a point]\label{def:c29-tensor}
On a $C^1$-manifold, a tensor of valence $(r,s)$ at $p$ is an element of
\[
  (T_pM)^r_s=\underbrace{T^*_pM\otimes\dots\otimes T^*_pM}_{r\text{ times}}\otimes\underbrace{T_pM\otimes\dots\otimes T_pM}_{s\text{ times}}.
\]
Its components in a chart $(U,\phi)\ni p$ are $t_\phi{}^{a_1\dots a_r}{}_{b_1\dots b_s}=t\bigl(d_p\phi^{a_1},\dots,d_p\phi^{a_r},\frac{\partial}{\partial\phi^{b_1}}\big|_p,\dots,\frac{\partial}{\partial\phi^{b_s}}\big|_p\bigr)$, and the tensor is recovered by
\[
  t=t_\phi{}^{a_1\dots a_r}{}_{b_1\dots b_s}\frac{\partial}{\partial\phi^{a_1}}\Big|_p\otimes\dots\otimes\frac{\partial}{\partial\phi^{a_r}}\Big|_p\otimes d_p\phi^{b_1}\otimes\dots\otimes d_p\phi^{b_s}.
\]
\end{definition}
\begin{recall}[Tensors]
\cref{def:c03-tensor}, \cref{def:c03-comp-tensor}.
\end{recall}
\begin{coursenote}[Which slot is which]
The paper writes $r$ factors of $T^*_pM$ and $s$ factors of $T_pM$ in $(T_pM)^r_s$, but then evaluates the \emph{upper} indices on the differentials $d_p\phi^a$ and the \emph{lower} ones on the basis vectors, and reads a vector as valence $(1,0)$ and a covector as $(0,1)$. So the printed factors are the other way round: valence $(r,s)$ means $r$ copies of $T_pM$ and $s$ of $T^*_pM$, as in chapter~\ref{ch:c03}. With this convention everything else in the paper is consistent.
\end{coursenote}

Remarks. (1) For $v\in V$, $w\in W$ the element $v\otimes w$ is the bilinear form on $V^*\times W^*$ given by $(v\otimes w)(\nu,\mu)=v(\nu)w(\mu)$; the space $V\otimes W=\{a^{ij}v_i\otimes w_j\}$ (bases $v_i$, $w_j$) is the vector space of such bilinear forms, with the usual sum and scalar multiple. (2) $(T_pM)^r_s$ is a vector space of dimension $(\dim M)^{r+s}$; a covector has valence $(0,1)$, a vector (by $V=(V^*)^*$) valence $(1,0)$, a scalar $(0,0)$. (3) Under a change of chart $(U,\phi)\to(V,\psi)$ the components change by
\[
  t_\psi{}^{a_1\dots a_r}{}_{b_1\dots b_s}=(C^\psi_\phi)^{a_1}{}_{m_1}\cdots(C^\psi_\phi)^{a_r}{}_{m_r}(C^\phi_\psi)^{n_1}{}_{b_1}\cdots(C^\phi_\psi)^{n_s}{}_{b_s}\,t_\phi{}^{m_1\dots m_r}{}_{n_1\dots n_s}
\]
(the printed formula has $t_\psi$ on the right-hand side as well; the components on the right are those of the old chart $\phi$).

\begin{definition}[Tensor bundle, tensor field]\label{def:c29-field}
The tensor bundle $(TM)^r_s=\{(p,t)\mid p\in M,\,t\in(T_pM)^r_s\}$ with the projection $\pi(p,t)=p$. A tensor field is a map $T:M\to(TM)^r_s$ with $\pi\circ T=\mathrm{id}_M$.
\end{definition}
\begin{recall}[Bundles and sections]
\cref{def:c06-bundle}, \cref{def:c06-section}, \cref{def:c06-vf}.
\end{recall}

Remarks. (1) A chart induces the vector fields $\frac{\partial}{\partial\phi^a}:U\to(TU)^1_0$, $p\mapsto(p,\frac{\partial}{\partial\phi^a}\big|_p)$ and the covector fields $d\phi^a:U\to(TU)^0_1$, $p\mapsto(p,d_p\phi^a)$. (2) A tensor field in terms of component functions is
\[
  T_\phi(p)=T_\phi{}^{a_1\dots a_r}{}_{b_1\dots b_s}(\phi(p))\frac{\partial}{\partial\phi^{a_1}}\otimes\dots\otimes\frac{\partial}{\partial\phi^{a_r}}\otimes d\phi^{b_1}\otimes\dots\otimes d\phi^{b_s},
\]
and $T$ is a $C^m$-tensor field if the component functions $\phi(U)\to\R$ are $m$ times continuously differentiable.

\section{Lecture III: Hyperbolic geometries}\label{sec:c29-L3}

\obg{Lecture III} The real work begins. The question is which geometries on a smooth manifold are candidates for the geometry of spacetime. The starting point is the loosest possible assumption, an arbitrary tensor field; the task is to find what it must satisfy. The definition of spacetime itself is only reached in Lecture V.

\begin{definition}[New concepts: Geometry]\label{def:c29-geometry}
A geometry on a smooth manifold $M$ is a smooth tensor field $G$, possibly restricted by further conditions.
\end{definition}

\begin{example}[Examples of geometries]\label{ex:c29-geometries}
\begin{enumerate}
  \item \emph{Riemannian geometry}: a $(0,2)$-tensor field $g$ that is symmetric, $g(X,Y)=g(Y,X)$, and positive definite, $g(X,X)>0$ for $X\ne0$. The geometry of general and special relativity is \emph{Lorentzian}: a symmetric $(0,2)$-tensor $g$ of signature $(1,\dim M-1)$.
  \item \emph{Symplectic geometry} of the phase space of mechanical systems: a $(0,2)$-tensor field $\omega$ that is antisymmetric, $\omega(X,Y)=-\omega(Y,X)$; non-degenerate ($\omega(X,Y)=0$ for all $X$ implies $Y=0$); and closed, $\partial_a\omega_{bc}+\partial_b\omega_{ca}+\partial_c\omega_{ab}=0$.
  \item \emph{Area metric geometry}, an example of spacetime geometry beyond the Lorentzian one: a $(0,4)$-tensor field $G$ with the symmetries
  \begin{gather*}
    G(X,Y,A,B)=G(A,B,X,Y),\\ G(A,B,X,Y)=-G(B,A,X,Y),
  \end{gather*}
  and non-degenerate: if $G(A,B,X,Y)=0$ for every linearly independent pair $A,B$, then $X,Y$ must be linearly dependent. To be a viable spacetime structure an area metric will need further algebraic conditions, which are exactly the conditions that give the Lorentzian signature in the metric case.
\end{enumerate}
\end{example}
\begin{coursenote}[Closedness as printed]
The paper writes the closedness condition with $\partial/\partial\phi^a$ in place of $\partial_a$; and in the non-degeneracy of $\omega$ the text ``$\omega(X,Y)$ for all vectors $X$'' lacks ``$=0$''; the intended statement is the one above.
\end{coursenote}

Remarks of the paper. (1) A geometry is \emph{flat} if there is an atlas in which all component functions of $G$ are constant; flatness is not assumed in these lectures. (2) For Riemannian geometry $(M,g)$ flatness is equivalent to the vanishing of the Riemann--Christoffel tensor of $g$ (chapter~\ref{ch:c08}). For a symplectic geometry, instead, there are always ``Darboux'' charts in which $\omega$ has constant components, so every symplectic manifold is flat.

\begin{definition}[New concepts: Probing matter]\label{def:c29-probing}
Probing matter on a smooth manifold $M$ with geometry $G$ is a tensor field $\Phi$ with values in a vector subspace $V$ of some $(TM)^r_s$, satisfying a \emph{linear} partial differential equation
\[
  \Bigl[\sum_{n=1}^{s}Q_{MN}^{i_1\dots i_n}(x)\,\partial_{i_1}\cdots\partial_{i_n}\Bigr]\Phi^N(x)=0,
\]
where the coefficients $Q$ are made from the components of $G$ and their partial derivatives alone, and are such that (a) the equation as a whole is tensorial and (b) the initial value problem is well-posed. Small Latin indices run over $1,\dots,\dim M$, capital ones over $1,\dots,\dim V$.
\end{definition}

Remarks. (1) The name ``probing'' refers to the linearity: the lessons one can learn about the geometry from a field obeying a linear equation are not entangled with properties of the matter field itself, as they would be for a non-linear equation. (2) A restriction of $\Phi$ to a proper subspace $V$ is done by linear conditions, for instance symmetry: a generic $(0,2)$-tensor has $(\dim M)^2$ components, a symmetric one lives in a subspace of dimension $\dim M(\dim M+1)/2$. (3) After a change of chart, the top-order coefficient $Q^{i_1\dots i_s}_{MN}$ is the only one that still behaves as a tensor in every index; lower-order coefficients generally acquire additional terms built from the higher ones. (4) Well-posedness of the initial value problem means that on some hypersurfaces the data fix $\Phi$ everywhere through the equation, forwards or backwards in time; this is what the paper regards as the core of classical physics.

\begin{coursenote}[The range of the sum]
The sum in the definition runs from $n=1$ to $s$, but the Klein--Gordon example of the same lecture has the term $m^2\Phi$ without derivatives, so the sum must start at $n=0$ (compare the student question of lecture 26, chapter~\ref{ch:c26}, \cref{def:c26-principal}). Also, the letter $s$ is used for the order of the equation and, in $(TM)^r_s$, for the valence: two different quantities.
\end{coursenote}

\begin{theorem}[Hyperbolicity of the principal polynomial]\label{thm:c29-hyp}
Suppose the linear equations just defined admit a well-posed initial value problem. Then
\[
  P(x,k)=\pm\rho\,\det_{M,N}\bigl[Q_{MN}^{i_1\dots i_s}(x)\,k_{i_1}\cdots k_{i_s}\bigr]
\]
is a positive hyperbolic homogeneous polynomial on every cotangent space $T^*_xM$; here $\rho$ is a scalar density made from $G$ which makes $P(x,k)$ a function on the cotangent bundle. The terms mean:
\begin{enumerate}
  \item \emph{homogeneous} means $P(x,\lambda k)=\lambda^{\deg P}P(x,k)$ for all real $\lambda$ and covectors $k\in T^*_xM$; $\deg P$ is the degree of $P$;
  \item \emph{hyperbolic} means that there is a covector $h$ (a hyperbolic covector of $P$ at $x$) with $P(x,h)\ne0$ such that for every $q\in T^*_xM$ the equation $P(x,q+\lambda h)=0$ in the real unknown $\lambda$ has $\deg P$ real solutions (and hence only real ones);
  \item \emph{positive} means that the sign $\pm$ is chosen so that $P(h)>0$.
\end{enumerate}
\end{theorem}
\begin{recall}[Principal polynomial and hyperbolicity in the lecture]
\cref{def:c26-principal}, \cref{def:c26-hyp}. In the hyperbolic condition the solutions are counted with algebraic multiplicity, as in the lecture. A line through the origin of a cone touches it once geometrically, but the root counts with its multiplicity.
\end{recall}
The paper states this theorem without proof (the underlying PDE theory is in H\"ormander, [3] of the paper).

\begin{nfigure}
\begin{tikzpicture}[scale=0.8]
  \begin{scope}
    \draw[->,gray] (-2.6,0)--(2.6,0) node[right,font=\scriptsize]{$k_1$};
    \draw[->,gray] (0,-2.2)--(0,2.4) node[above,font=\scriptsize]{$k_0$};
    \draw[thmcol,very thick] (-2.2,-2.2)--(2.2,2.2);
    \draw[thmcol,very thick] (-2.2,2.2)--(2.2,-2.2);
    \draw[defcol,thick,dashed] (1.5,-2.1)--(1.5,2.3);
    \fill[defcol] (1.5,1.5) circle (2pt) (1.5,-1.5) circle (2pt);
    \node[font=\scriptsize,text=defcol] at (2.0,2.45) {line $q+\lambda h$};
    \node[font=\scriptsize] at (0.0,-2.7) {degree 2: $g^{ab}k_ak_b=0$, 2 intersections};
  \end{scope}
  \begin{scope}[xshift=8.2cm]
    \draw[->,gray] (-2.6,0)--(2.6,0) node[right,font=\scriptsize]{$k_1$};
    \draw[->,gray] (0,-2.2)--(0,2.4) node[above,font=\scriptsize]{$k_0$};
    \foreach \m in {1,2}{
      \draw[thmcol,very thick] (-2.2/\m,-2.2)--(2.2/\m,2.2);
      \draw[thmcol,very thick] (-2.2/\m,2.2)--(2.2/\m,-2.2);}
    \draw[defcol,thick,dashed] (0.9,-2.1)--(0.9,2.3);
    \fill[defcol] (0.9,0.9) circle (2pt) (0.9,-0.9) circle (2pt) (0.9,1.8) circle (2pt) (0.9,-1.8) circle (2pt);
    \node[font=\scriptsize] at (0.0,-2.7) {degree 4: two nested cones, 4 intersections};
  \end{scope}
\end{tikzpicture}
\end{nfigure}
\begin{coursenote}[Figure, drawn here]
Instead of the paper's figure, a drawing in 1+1 dimensions of the null set of a polynomial of degree two ($k_0^2-k_1^2$) and degree four ($(k_0^2-k_1^2)(k_0^2-4k_1^2)$, my own example of a hyperbolic quartic) in the cotangent plane: any line parallel to $h=(1,0)$ (the $k_0$ direction) meets the null set in exactly $\deg P$ points.
\end{coursenote}

Remarks. (1) Rescaling $P$ by a scalar density $\rho$ changes neither hyperbolicity nor any other condition revealed in these lectures. (2) Only a hypersurface $\Sigma$ whose \emph{conormals} (covectors $n\in T^*_yM$ with $n(\dot\gamma_y)=0$ for all smooth curves $\gamma$ in $\Sigma$ through $y$) are hyperbolic covectors can be an initial data surface of the field equations: hyperbolicity of $P$ is the statement that there are viable initial data surfaces at all. (3) That $P$ be hyperbolic restricts the geometry $G$; finding the extent of this restriction is one main topic of the lectures.

\begin{example}[Klein--Gordon on a Lorentzian manifold]\label{ex:c29-KG}
For $(M,g)$ the Klein--Gordon equation
\[
  g^{ab}(x)\partial_a\partial_b\Phi-\tfrac12g^{ab}g^{ms}\bigl(\partial_ag_{sb}+\partial_bg_{as}-\partial_sg_{ab}\bigr)\partial_m\Phi+m^2\Phi=0
\]
has $P(x,k)=g^{ab}(x)k_ak_b$, which is hyperbolic if and only if $g$ has signature $(1,d-1)$ or $(d-1,1)$; positivity selects $(1,d-1)$. Hence Klein--Gordon matter can be predictive only on a Lorentzian geometry. The hypersurfaces with hyperbolic conormal are the spacelike ones, and indeed every initial data surface of Klein--Gordon is spacelike (the converse is not true). The theorem generalises this to any tensorial geometry.
\end{example}
\begin{coursenote}[Sign of the mass term]
The paper writes $+m^2\Phi$ here and $-m^2\Phi$ in Lecture IV; the sign is irrelevant for $P$, which depends only on the highest order. Also the first order coefficient in the first-derivative term is the Christoffel symbol contraction $g^{ab}\Gamma^m{}_{ab}$ up to sign: it is not a tensor, as Remark 3 above says.
\end{coursenote}

\begin{definition}[Hyperbolicity cone]\label{def:c29-cone}
If $h$ is a hyperbolic covector of $P$, there is an entire connected set of hyperbolic covectors around it, the hyperbolicity cone $C(P,h)$. It is constructed as follows. Expand
\[
  P(x,q+\lambda h)=h_0(x,q,h)\lambda^{\deg P}+h_1(x,q,h)\lambda^{\deg P-1}+\dots+h_{\deg P}(x,q,h),
\]
and build the matrices
\[
  H_i(q,h)=\begin{bmatrix}
    h_1&h_3&h_5&\dots&h_{2i-1}\\
    h_0&h_2&h_4&\dots&h_{2i-2}\\
    0&h_1&h_3&\dots&h_{2i-3}\\
    0&h_0&h_2&\dots&h_{2i-4}\\
    \vdots&\vdots&\vdots&&\vdots\\
    0&0&0&\dots&h_i
  \end{bmatrix}_{i\times i},\qquad h_j\to0\text{ for }j>\deg P\ [\text{the paper writes }j>i],
\]
for $i=1,\dots,\deg P$. Then
\[
  C(P,h)=\{h'\in T^*_xM\mid\det H_i(h',h)>0\text{ for all }i=1,\dots,\deg P\}.
\]
\end{definition}
\begin{coursenote}[Hurwitz matrices]
The matrices have the Hurwitz (Routh--Hurwitz) shape. The paper says ``$h_j$ is set to $0$ for $j>i$''; the coefficients exist only up to $h_{\deg P}$, so the zero entries are those with index above $\deg P$ (in the example below only $h_3=0$ appears for $\deg P=2$).
\end{coursenote}

Remarks. (1) By construction every $h'\in C(P,h)$ gives the same cone, $C(P,h')=C(P,h)$, and the cone is open in the standard topology of the cotangent space. G{\aa}rding showed that it is a \emph{convex cone}: $\lambda h'\in C(P,h)$ and $h'+h''\in C(P,h)$ for $h',h''\in C(P,h)$ and positive real $\lambda$ (stated without proof, compare \cref{thm:c26-garding}). (2) $P$ is positive on the whole cone and vanishes on its boundary $\partial C$. (3) If $P$ factorizes, $P=P_1^{\alpha_1}\cdots P_f^{\alpha_f}$, then $P$ is hyperbolic in the direction $h$ exactly when each $P_i$ is. For technical reasons, unless all $\alpha_i=1$ one replaces $P$ at each point by the \emph{reduced polynomial} $P=P_1\cdots P_f$. In any case
\[
  C(P,h)=C(P_1,h)\cap\dots\cap C(P_f,h).
\]

\begin{example}[Hyperbolicity cone of the Lorentzian polynomial]\label{ex:c29-lorcone}
For $P(k)=g^{-1}(k,k)$ with $g$ of signature $(+-\dots-)$, every $h$ with $g^{-1}(h,h)>0$ is hyperbolic. For such an $h$,
\[
  P(q+\lambda h)=\lambda^2g^{-1}(h,h)+2\lambda g^{-1}(q,h)+g^{-1}(q,q),
\]
so $h_0=g^{-1}(h,h)$, $h_1=2g^{-1}(q,h)$, $h_2=g^{-1}(q,q)$. Hence $\det H_1=2g^{-1}(q,h)$ and
\begin{gather*}
  H_2(q,h)=\begin{bmatrix}2g^{-1}(q,h)&0\\g^{-1}(h,h)&g^{-1}(q,q)\end{bmatrix},\\ \det H_2=2g^{-1}(q,h)\,g^{-1}(q,q)
\end{gather*}
(the paper prints the matrix after the word ``$\det$'' and does not evaluate it). Therefore the hyperbolicity cone containing $h$ is
\[
  C(P,h)=\{h'\in T^*_xM\mid g^{-1}(h',h')>0\text{ and }g^{-1}(h',h)>0\},
\]
the set of timelike covectors in the same half as $h$ (the paper writes $q$ in place of $h'$ in the second condition).
\end{example}
\begin{recall}[The light cone]
For a Lorentzian metric this is the open inside of the future (or past) light cone in the cotangent space, chapters~\ref{ch:c10} and~\ref{ch:c13}.
\end{recall}

\begin{exercise}\label{ex:c29-quartic}
Take $P(k)=(k_0^2-k_1^2)(k_0^2-4k_1^2)$ on $\R^2$ and $h=(1,0)$. Check by hand that $P$ is hyperbolic with respect to $h$ (four real roots $\lambda$ for every $q$) and find $C(P,h)$.
\end{exercise}

\section{Lecture IV: Massless dispersion}\label{sec:c29-L4}

\obg{Lecture IV} The hyperbolic polynomial gives the dispersion relation of the matter equations in the high-energy (equivalently, massless) limit. This is used to derive the action of a massless point particle in terms of a \emph{dual polynomial}, which then enters the definition of spacetime.

\begin{definition}[New concepts: Locally wavelike solutions]\label{def:c29-wavelike}
A family of locally wavelike solutions of linear matter equations is a family of solutions
\[
  \Phi^N_\lambda(x)=A^N_\lambda(x)\,e^{iS(x)/\lambda}\qquad\text{for }\phi(x)\in U,
\]
where $(U,\phi)$ is a conveniently small chart, $\lambda>0$ is a parameter, the \emph{phase function} $S$ is smooth, real and has nowhere vanishing differential, and the real amplitude components have Taylor expansions $A^N_\lambda=\sum_{n=0}^\infty a^N_n(x)\lambda^n$ with $a_0(x)\ne0$.
\end{definition}

Remarks. (1) A hypersurface $S=\text{const}$ is a \emph{wave front}, $d_pS/\lambda$ is the \emph{wave covector} at $p$. Later a \emph{ray vector} is associated with each wave covector. (2) Inserting the ansatz in the field equations gives
\[
  Q_{MN}^{i_1\dots i_s}(x)\,\partial_{i_1}\cdots\partial_{i_s}S(x)\,a^N_0(x)+(\dots)_M\lambda+(\dots)_M\lambda^2+\dots=0,
\]
so to lowest order one needs $\det_{M,N}\bigl(Q_{MN}^{i_1\dots i_s}\partial_{i_1}S\cdots\partial_{i_s}S\bigr)=0$. In terms of the function $P$ of the previous lecture the wave covector $k=dS/\lambda$ must, to lowest order, satisfy the \emph{dispersion relation}
\[
  P(x,k)=0.
\]
(3) Taking into account higher orders in $\lambda$, the dispersion relation may or may not be modified in the high-frequency limit $\lambda\to0$. If it is modified the field equations are called \emph{massive}, otherwise \emph{massless}.

\begin{example}[Klein--Gordon in flat spacetime]\label{ex:c29-KGflat}
In a suitable chart the flat Klein--Gordon equation is $\eta^{ab}\partial_a\partial_b\Phi-m^2\Phi=0$, $m\ge0$. Inserting the ansatz gives
\begin{multline*}
  -\eta^{ab}\partial_aS\partial_bS\,A_\lambda+\bigl[2i\eta^{ab}\partial_aS\partial_bA_\lambda+i\eta^{ab}\partial_a\partial_bA_\lambda\bigr]\lambda\\+\bigl[\eta^{ab}\partial_a\partial_bA_\lambda+m^2A_\lambda\bigr]\lambda^2=0.
\end{multline*}
So to lowest order $S(x)=k_ax^a$ solves it for any covector with $\eta^{ab}k_ak_b=0$, and the equation reduces to
\begin{gather*}
  \sum_{j=0}^\infty\bigl[\eta^{ab}\partial_a\partial_ba_{j-1}+m^2a_{j-1}+2I\eta^{ab}k_a\partial_ba_j\bigr]\lambda^{j-1},\\ a_{-1}=0,
\end{gather*}
solved for arbitrary $\lambda$ by
\[
  a_n(x)=a\frac1{n!}\Bigl(\frac{i\,m^2l_cx^c}{2\eta^{ab}k_al_b}\Bigr)^n,\qquad A_\lambda(x)=a\,e^{i\frac{\lambda m^2}{2\eta^{ab}l_ak_b}l_cx^c},
\]
for a real constant $a$ and a covector $l$ with $\eta^{ab}l_al_b=0$, $\eta^{ab}l_ak_b\ne0$. To keep the amplitude real, absorb the phase into $S$:
\[
  S(x)=\Bigl(1+\frac{\lambda^2m^2}{2\eta^{ab}l_ak_b}\Bigr)k_cx^c,\qquad A(x)=a.
\]
Taking all orders of $\lambda$ into account one gets the dispersion relation $P(x,\partial S(x))=m^2$. So Klein--Gordon with $m=0$ is massless, with $m>0$ massive \obg{Lecture IV, Example}.
\end{example}
\begin{coursenote}[Slips in the printed example]
The printed series has $I$ where the imaginary unit $i$ is meant. The printed expressions for $a_n$ and for $S$ were not rederived here; they are given as printed [?].
\end{coursenote}

\begin{definition}[New concepts: Massless momenta]\label{def:c29-null}
The set of massless momenta at $x\in M$ is the cone
\[
  N_x=\{k\in T^*_xM\mid P(x,k)=0\}.
\]
\end{definition}

Remarks. (1) For technical precision one often uses the \emph{smooth subcone} $N_x^{\rm smooth}=\{k\in N_x\mid DP(x,k)\ne0\}$, where $DP$ is the derivative of the reduced $P$ along the cotangent fibre. (2) The cone $N_x$ is determined by $P(x,\cdot)$; conversely, it determines $P_x$ up to a constant factor. This is subtle but central, and follows from the real Nullstellensatz
\[
  \mathcal I(N_x)=\{\alpha P_x\mid\alpha\in\R\},
\]
valid for any reduced hyperbolic polynomial. Here $\mathcal I(S)$ is the set of all polynomials on $T^*_xM$ vanishing on $S\subset T^*_xM$. The Nullstellensatz is used shortly.

To attach velocities to massless momenta in a physically meaningful way, one uses the dynamics of free massless point particles.

\begin{theorem}[Massless point particle]\label{thm:c29-massless}
The action of a free massless point particle is
\[
  I_0[x,q,\lambda]=\int d\tau\,\bigl[q_a\dot x^a+\lambda P(x,q)\bigr],
\]
where $\lambda$ is a Lagrange multiplier function.
\end{theorem}

Remarks. (1) By construction this describes a free massless particle; the geometry enters only through the dispersion relation enforced by the Lagrange multiplier term. (2) To eliminate $q$ and $\lambda$ and obtain an action for the trajectory $x$ alone: variation with respect to $\lambda$ enforces the null condition, variation with respect to $q$ gives $\dot x=\lambda\,DP_x(q)$ for $q\in N^{\rm smooth}$, which implies the weaker equation $[DP_x(q)]=[\dot x/\lambda]$, where $[X]$ is the projective class of all vectors collinear with $X$. This latter equation must be inverted.

\begin{definition}[Dual polynomial]\label{def:c29-dual}
The polynomial $P^\#_x:T_xM\to\R$ is \emph{dual} to an \emph{irreducible} polynomial $P:T^*_xM\to\R$ if
\[
  P^\#_x\bigl(DP_x(N^{\rm smooth}_x)\bigr)=0,
\]
which determines it up to real scale. For a reducible $P$ with irreducible factors $P_i$ one defines the dual as the product
\[
  (P_1(x,k)\cdots P_f(x,k))^\#(x,v)=P_1^\#(x,v)\cdots P_f^\#(x,v);
\]
so $P^\#$ is determined up to real scale and again satisfies the equation above.
\end{definition}
\begin{recall}[Dual polynomial in the lecture]
\cref{def:c26-dual}; there it was introduced geometrically as the polynomial vanishing on the image of the null surface under the fibre derivative.
\end{recall}

Remarks. (1) Geometrically, $P^\#_x$ is the polynomial whose zero set is the \emph{dual null cone} $N^\#_x$, the image of $N_x$ under the gradient map $DP:N_x\to T_xM$, $k\mapsto\partial P/\partial k_a(x,k)$. (2) The existence of $P^\#_x$, and its computability by algorithm for any reduced hyperbolic $P$, hinges on the real Nullstellensatz; the degree of $P^\#$ is in general different from that of $P$. (3) In principle $P^\#$ is computed with Buchberger's algorithm (elimination theory); in practice, for high degree and several variables, this exhausts current computer algebra systems. In some physically interesting cases one can guess $P^\#$ from physical reasoning and check it afterwards. Since a dual always exists for reduced hyperbolic $P$, the lectures simply assume one has been found.

\begin{theorem}[Gauss map and dual Gauss map]\label{thm:c29-gauss}
The \emph{Gauss map} $[DP]:[N^{\rm smooth}_x]\to[N^\#_x]$, $[q]\mapsto[DP_x(q)]$, and the \emph{dual Gauss map} $[DP^\#]:[N_x^{\#,\rm smooth}]\to[N_x]$, $[X]\mapsto[DP^\#_x(X)]$, are partial inverses of each other: for $k\in N^{\rm smooth}_x$
\[
  [DP^\#_x]\bigl([DP_x]([k])\bigr)=[k]\qquad\text{if }\det(DDP_x)(k)\ne0,
\]
and likewise with $P$ and $P^\#$ exchanged.
\end{theorem}
\begin{proof}
Write the defining property of the dual in the form (which needs no explicit restriction to null covectors)
\[
  P^\#(x,DP(x,k))=Q(k)\,P(k)\qquad\text{for all covectors }k.
\]
Differentiate with respect to $k_a$, using the chain rule on the left: $DP^\#_b(DP(k))\,\partial_a\partial^bP(k)$. On the right, the term with $\partial_aQ$ is multiplied by $P(k)$, which vanishes for null $k$, so for null $k$ one is left with $Q(k)\,\partial_aP(k)$. Now apply Euler's theorem to the function $DP$, which is homogeneous of degree $\deg P-1$: $DDP(k)\,k=(\deg P-1)\,DP(k)$ (this is the derivative form of $DP(k)\,k=(\deg P)P(k)$). Hence the vector $\frac{Q}{\deg P-1}k$ solves the linear equation $DDP(k)\,v=Q\,DP(k)$ which $DP^\#(DP(k))$ solves, and if $\det(DDP)(k)\ne0$ the solution is unique:
\[
  DP^\#(x,DP(x,k))=\frac{Q(x,k)}{\deg P-1}\,k.
\]
In projective language this is the statement of the theorem.
\end{proof}
\begin{coursenote}[What was added to the proof]
The paper gives the proof in two lines (``differentiation, chain rule, Euler''). The intermediate step (the term with $\partial Q$ vanishing on the null set, and the use of Euler's identity for $DP$ to identify the solution of the linear system) is spelled out here; the result is the one printed.
\end{coursenote}

Remarks. (1) Thus the momentum-velocity relation of the massless particle can be solved, up to scale, for the momentum: $[q]=[DP^\#_x]([\dot x/\lambda])$, and by homogeneity of $DP^\#_x$ and the projective brackets the function $\lambda$ can be dropped. Without brackets another undetermined function $\mu$ appears,
\[
  q=\mu\,DP^\#_x(\dot x).
\]
(2) So the Gauss map $[DP_x]$ and its inverse $[DP^\#_x]$ pair null momenta in $N^{\rm smooth}_x$ with null velocities in $N^{\#,\rm smooth}_x$, up to a harmless conformal factor. (3) Insert $q$ into the action and use Euler's theorem again, for the homogeneous polynomial $P^\#_x$: one finds the massless point particle action
\[
  I_0[x,\mu]=\int d\tau\,\mu\,P^\#(x,\dot x).
\]
The Lagrange multiplier $\mu$ appears again, as it must to enforce the null constraint $P^\#_x(\dot x)=0$. In this way $P^\#$ acquires a direct physical role: it is the geometry on tangent space that massless particles feel.
\begin{recall}[Gauss map in the lecture]
\cref{def:c26-maps}, and the physical meaning in \cref{sec:c26-meaning}.
\end{recall}

\section{Lecture V: Spacetimes}\label{sec:c29-L5}

\obg{Lecture V} Now a definition of a spacetime geometry. The requirements are needed so that the probing matter is predictive and so that observers do not disagree about which particles are particles and which are antiparticles; they are also enough to rebuild the kinematics of general relativity on any geometry that satisfies them.

\begin{definition}[Bi-hyperbolic polynomial]\label{def:c29-bihyp}
$P_x$ is \emph{bi-hyperbolic} if both $P_x$ and its dual $P^\#_x$ are hyperbolic.
\end{definition}

Remarks. (1) Hyperbolicity of a polynomial on tangent space is defined as on cotangent space, with covectors replaced by vectors. (2) Hyperbolic does not imply bi-hyperbolic. A notable exception: degree two (necessarily an inverse Lorentzian metric, $P(k)=g^{-1}(k,k)$): its dual $P^\#(X)=g(X,X)$ is hyperbolic if and only if $P$ is, because $g$ has Lorentzian signature if and only if $g^{-1}$ has.
\begin{recall}[Bihyperbolicity in the lecture]
\cref{thm:c26-bihyp}, there stated in words with the physical reading: predictivity needs $P$ hyperbolic, quantizability needs $P^\#$ hyperbolic.
\end{recall}

\begin{definition}[New concepts: Time orientation, observer, energy]\label{def:c29-timeor}
Take a smooth vector field $T$ without zeros such that $T(x)$ lies, at each $x\in M$, inside a hyperbolicity cone of $P^\#_x$; call that cone $C^\#_x$. The smooth distribution of cones $C^\#=\{(x,C^\#_x)\mid x\in M\}$ is a \emph{time orientation} of $(M,G)$.
\end{definition}

Remarks. (1) A time orientation is needed for a meaningful definition of observers: the worldline $x:\R\to M$ of an \emph{observer} must have tangent vectors $\dot x_{x(\tau)}\in C^\#_{x(\tau)}$ for all $\tau$, i.e.\ the observer travels into the future directions of $C^\#$. (2) A full definition of observers, including their purely spatial directions, is postponed to the next lectures (Lecture VI). (3) Already with the partial definition one can define the \emph{energy} of a covector $q\in T^*_xM$ seen by an observer with tangent vector $X\in C^\#$ as $q(X)$; it only has physical meaning for covectors that are particle momenta, treated later.

\begin{definition}[New concepts: Cone of positive energy covectors]\label{def:c29-poscone}
At $x\in M$,
\[
  (C^\#_x)^+=\{p\in T^*_xM\mid p(X)>0\text{ for all }X\in C^\#_x\}.
\]
\end{definition}

Remarks. (1) For a covector $q\in(C^\#_x)^+$ the energy $q(X)$ is positive for every observer $X\in C^\#_x$ at that point; this is the reason the cone $C^\#_x$ is defined at all. (2) If the reduced $P$ factorises, $P=P_1\cdots P_f$, then the positive energy cone is the sum of those of the factors,
\[
  (C^\#)^+=(C^\#_1)^++\dots+(C^\#_f)^+,
\]
where the sum of two subsets of a vector space is the set of all sums of one element of each.

\begin{definition}[New concepts: Energy-distinguishing polynomial]\label{def:c29-endist}
A bi-hyperbolic $P_x$ is \emph{energy-distinguishing} if every massless momentum has a definite energy sign, that is, if
\[
  N_x=N^+_x\,\dot\cup\,N^-_x,\qquad N^+_x=N_x\cap(C^\#)^+,\quad N^-_x=N_x\cap-(C^\#)^+ .
\]
\end{definition}

\begin{remark}[Consequences of energy distinction]\label{rem:c29-evenness}
(1) If the principal polynomial is bi-hyperbolic and energy-distinguishing at every point, the sign of the energy of a massless momentum is observer independent; on a flat $G$ this gives a unique split of fundamental solutions into positive and negative energy parts. (2) Energy-distinguishing hyperbolic polynomials have \emph{even degree}.
\end{remark}
\begin{proof}
First one proves that bi-hyperbolicity implies
\[
  \operatorname{closure}\bigl((C^\#_x)^+\bigr)\cap-\operatorname{closure}\bigl((C^\#_x)^+\bigr)=\{0\}.
\]
Let $k_0$ lie in both sets. By the definition of the dual cone, $X\cdot k_0\ge0$ and $X\cdot k_0\le0$ for all $X\in C^\#_x$, i.e.\ $X\cdot k_0=0$. Then the hyperbolicity cone $C^\#_x$ would lie in a plane (if $k_0\neq0$), contradicting its openness. Second, assume $N_x$ contains a plane. The first step puts $(C^\#_x)^+\setminus\{0\}$ strictly inside a half-space, which cannot hold a whole plane through the origin; thus a null plane would rule out the energy-distinguishing property. Third, this restricts $P$ to even degree. Suppose $\deg P$ were odd: then there would be an odd number of null sheets; but by homogeneity null sheets in a cotangent space come in pairs, each the point reflection of its partner, so at least one null hyperplane would be present, contradicting the second step.
\end{proof}
\begin{coursenote}[Notation in the proof]
The paper uses the letter $x$ both for the point of $M$ and for the vectors in the dual-cone argument; the vectors are $X\in C^\#_x$. The steps are only sketched in the paper, as above.
\end{coursenote}

\begin{definition}[Spacetime with respect to matter dynamics]\label{def:c29-spacetime}
A geometry $(M,G,C^\#)$ is a \emph{spacetime with respect to matter dynamics} $D_{MN}(\partial)\Phi^N=0$ if the cotangent bundle function $P:T^*_xM\to\R$ defined by the principal polynomial of that linear matter dynamics is
\begin{enumerate}
  \item everywhere bi-hyperbolic,
  \item such that $C^\#$ is a time orientation defined in terms of $P$,
  \item everywhere energy-distinguishing with respect to this time orientation.
\end{enumerate}
\end{definition}
\begin{recall}[The refined definition of spacetime]
\cref{def:c26-spacetime}: manifold, geometry and matter action.
\end{recall}

Remarks. (1) Whether $(M,G)$ is a viable spacetime depends on the linear matter equations used to probe it; the paper counts this as an insight and not as a defect. (2) The restriction to linear matter makes the principal polynomial depend on $G$ only, not on solutions. A non-linear equation can be linearised about an exact solution, $\Phi=\Phi_{\rm exact}+\delta\Phi$; the principal symbol of the linearised equation still decides causality, but it now depends on the background solution $\Phi_{\rm exact}$, so the statements are no longer purely about the geometry. Linear matter is the right probe. (3) Lorentzian geometry $(M,g)$ is a spacetime with respect to Maxwell theory, as the next example shows.

\begin{example}[Maxwell theory on a Lorentzian manifold]\label{ex:c29-maxwell}
In Lorentz gauge the Maxwell equations read
\[
  \frac{\delta^n_m}{\sqrt{-\det g(x)}}\,\partial_a\Bigl(\sqrt{-\det g(x)}\,g^{ab}(x)\,\partial_bA_n(x)\Bigr)=0,
\]
and one reads off at each point $P_x(k)=g^{ab}(x)k_ak_b$, hyperbolic because $g^{-1}$ has Lorentzian signature. The polynomial $P^\#(X):=g_{ab}X^aX^b$ is dual to it, since for $k$ with $P(k)=0$
\[
  P^\#(DP(k))=g_{ab}(2g^{am}k_m)(2g^{bn}k_n)=4\,g^{mn}k_mk_n=0,
\]
and it is hyperbolic since $g$ is Lorentzian. Thus $P$ is bi-hyperbolic. It is energy-distinguishing: choose one of the two hyperbolicity cones of $P^\#$ as time orientation $C^\#$; then
\[
  (C^\#)^+=\{q\in T^*_xM\mid q(X)>0\text{ for all }X\in C^\#\}=\operatorname{closure}(C),
\]
where $C$ is the hyperbolicity cone of $P$ lying entirely within $(C^\#)^+$ (see Lecture VI). But then $N=N^+\,\dot\cup\,N^-$.
\end{example}
\begin{coursenote}[Corrections to the printed example]
The paper prints $g^{mn}k_nk_n$ and omits the factor 4; the correct value is $4g^{mn}k_mk_n$, which vanishes on null $k$. It also writes ``b-hyperbolic'' for bi-hyperbolic. In this Maxwell example the equations are those for $A_n$ in Lorentz gauge: the gauge freedom (where the determinant construction of the principal polynomial ``works differently'', chapter~\ref{ch:c26}) is hidden in this choice of gauge.
\end{coursenote}

\section{Lecture VI: Massive dispersion}\label{sec:c29-L6}

\obg{Lecture VI} Massive point particles are needed to complete the picture, and because they turn massive covectors into velocity vectors; a corollary singles out the physically acceptable modified dispersion relations.

\begin{definition}[New concepts: Massive positive energy momentum, mass]\label{def:c29-massive}
Fix a hyperbolicity cone $C\subset(C^\#)^+$. A covector $q\in C$ is called a \emph{massive positive energy momentum}, and its \emph{mass} is the number $m>0$ with $P(q)=m^{\deg P}$.
\end{definition}

Remarks. (1) Some hyperbolicity cone of $P$ always lies inside $(C^\#)^+$. Indeed the boundary of a hyperbolicity cone is a connected set of null covectors, while the null covectors split into the two disjoint pieces $(C^\#)^+$ and $-(C^\#)^+$; so $C$ or $-C$ lies in $(C^\#)^+$. (2) The positivity of $P$ on $C$ is guaranteed by the choice of overall sign in Lecture III.

\begin{theorem}[Barrier function and Legendre map]\label{thm:c29-legendre}
For a bi-hyperbolic and energy-distinguishing $P$ the \emph{barrier function}
\[
  f_x:C_x\to\R,\qquad f_x(q)=-\frac1{\deg P}\ln P(x,q)
\]
is strictly convex and essentially smooth; this makes the \emph{Legendre map}
\[
  L_x:C_x\to L_x(C_x)\subset T_xM,\qquad L_x(q):=-D(\ln P)(x,q)
\]
invertible, with
\[
  L_x^{-1}:L_x(C_x)\to C_x,\qquad L_x^{-1}(q)=-Df^L_x,
\]
given in terms of the \emph{Legendre transform} of the barrier function,
\[
  f^L_x:L_x(C_x)\to\R,\qquad f^L_x(X)=-L_x^{-1}(X)\,X-f_x\bigl(L_x^{-1}(X)\bigr),
\]
at every point $x\in M$.
\end{theorem}
\begin{recall}[Legendre map in the lecture]
\cref{def:c26-maps}, \cref{thm:c26-legendre}.
\end{recall}
The paper states this theorem without proof; the convex analysis is in Rockafellar [10].

Remarks. (1) \emph{Essential smoothness} concerns the behaviour of the barrier function near the boundary of the convex set where it is defined; see a text on convex analysis. The functions Legendre transformed in classical mechanics or thermodynamics live on a whole vector space, which is convex but has no boundary, so the condition is trivially satisfied and thus less known. (2) It takes bi-hyperbolicity and the energy-distinguishing property together to get an invertible Legendre map; all the conditions of Lecture V are thus used, and the map is put to work for massive particles next.

\begin{theorem}[Free massive particle]\label{thm:c29-massivepart}
A free particle of mass $m>0$ and positive energy moves according to the action
\[
  S[x]=m\int d\tau\,P^*_{x(\tau)}\bigl(\dot x(\tau)\bigr)^{1/\deg P},
\]
where on the whole cone $L_x(C_x)$ the function $P^*_x(X)=P_x\bigl(L_x^{-1}(X)\bigr)^{-1}$ is used.
\end{theorem}
\begin{proof}
The Helmholtz action
\[
  S[x,q,\lambda]=\int d\tau\Bigl[q_a\dot x^a-\lambda\,m\ln P\bigl(x,\tfrac qm\bigr)\Bigr]
\]
gives, when $\lambda$ is varied, the mass-shell condition for $q$; with the geometry-independent term $q_a\dot x^a$ it describes the particle, as in the massless case of Lecture IV. Varying $q$ gives $\dot x^a=(\lambda\deg P)\,L_x(q/m)$, which can be inverted by the Legendre theorem:
\[
  q=m\,L_x^{-1}\Bigl(\frac{\dot x}{\lambda\deg P}\Bigr).
\]
With this relation between massive momentum and the tangent vector to the worldline and with the definitions of barrier function and Legendre dual one eliminates $q$ and gets
\begin{multline*}
  S[x,\lambda]=-m\deg P\int d\tau\,\lambda f^L\Bigl(\frac{\dot x}{\lambda\deg P}\Bigr)\\=-m\deg P\int d\tau\,\bigl[\lambda f^L_x(\dot x)+\lambda\ln(\lambda\deg P)\bigr],
\end{multline*}
where the second equality uses the elementary scaling property $f^L(\alpha\dot x)=f^L(\dot x)-\ln\alpha$. Varying this action with respect to $\lambda$ gives
\[
  f^L(\dot x)+\ln(\lambda\deg P)+1=0.
\]
Using it twice, $\lambda f^L_x(\dot x)+\lambda\ln(\lambda\deg P)=-\lambda=-\exp\bigl(-f^L_x(\dot x)-1\bigr)/\deg P$. Since $\dot x\in L_x(C_x)$ one also has $L^{-1}(x,\dot x)(\dot x)=1$ and thus $f^L_x(\dot x)=-1-f_x\bigl(L^{-1}(\dot x)\bigr)$, and removing $\lambda$ leaves the equivalent action
\[
  S[x]=m\int d\tau\,P^*(x,\dot x)^{1/\deg P},
\]
which is the claim.
\end{proof}
\begin{coursenote}[Normalisations in this proof]
The proof is reproduced as printed. The normalisation of the Legendre map is not the same in all the formulas: with $L=-D\ln P$ and Euler's theorem one has $L_x(q)(q)=-\deg P$, whereas the proof uses $L^{-1}(\dot x)(\dot x)=1$ and $\dot x=(\lambda\deg P)L_x(q/m)$ (from $\dot x^a=\lambda\,\partial^a\ln P$ one would get $-\lambda L$). The structure of the argument, which is what matters, is: eliminate $q$ with the inverse Legendre map, then eliminate the multiplier $\lambda$ with its own equation, and use the homogeneity of $P^*$; I did not try to reconcile the constants [?]. The paper writes $L_x^{-1}(q)=-Df_x^L$ with $q$ where $X$ is meant.
\end{coursenote}

Remarks. (1) The tangent bundle function $P^*$ is in general not a polynomial, but it is easily seen to be homogeneous of degree $\deg P$ (the map $L^{-1}$ is homogeneous of degree $-1$ and $P$ of degree $\deg P$). (2) The action is reparametrisation invariant, as it should be. Parametrisations with $P(x,L^{-1}(x,\dot x))=1$ along the curve are distinguished: for them
\[
  \dot x=L_x(q/m)
\]
holds everywhere along the trajectory. One chooses such clocks and calls their time \emph{proper time}. (3) This establishes the physical meaning of the Legendre map: it maps massive positive energy momenta to the tangent vectors of their worldlines. The open convex cone $L_x(C_x)$ is therefore the \emph{cone of massive particle velocities}, and $P^*$ is the \emph{massive dual} of $P$, the tangent bundle geometry seen by massive particles. (4) It can be proved that $C^\#_x\subseteq L_x(C_x)$, so observers may be taken to be massive. The converse is false, the inclusion is generically proper.

\begin{definition}[New concepts: Observer]\label{def:c29-observer}
An observer is a curve $e:\R\to LM$ into the frame bundle of $M$ such that, along it,
\begin{enumerate}
  \item the first frame vector $e_0$ at each point is the tangent vector to the projected curve $\pi\circ e:\R\to M$,
  \item $e_0$ lies in the observer cone $C^\#$ at each point of the curve,
  \item $L^{-1}(e_0)(e_\alpha)=0$ for the remaining frame vectors $e_\alpha$, $\alpha=1,\dots,\dim M-1$.
\end{enumerate}
\end{definition}

Remarks. (1) The observer's spatial directions form the subspace $V_x=\{X\in T_xM\mid L^{-1}(e_0)(X)=0\}$. (2) An observer frame induces a unique dual frame $\epsilon^0=L^{-1}(e_0),\epsilon^1,\dots,\epsilon^{\dim M-1}$, and the zero component $q_0$ of $q=q_0\epsilon^0+q_\alpha\epsilon^\alpha$ is exactly the energy of the covector as seen by the observer.

\begin{theorem}[Dispersion relation of causal matter]\label{thm:c29-dispersion}
Matter of mass $m>0$ that propagates causally and whose energy sign all observers agree on obeys a dispersion relation of the form
\[
  P\bigl(E\,\epsilon^0+p_\alpha\epsilon^\alpha\bigr)=m^2
\]
with $P$ bi-hyperbolic and energy-distinguishing and $\epsilon^0,\dots,\epsilon^{\dim M-1}$ an observer coframe. Solving for $E$ relates the observer's energy $E$ to the spatial part $p_\alpha\epsilon^\alpha$ of the momentum $p\in C$.
\end{theorem}
\begin{coursenote}[Power of the mass]
With the definition of mass above, $P(q)=m^{\deg P}$, the right-hand side should be $m^{\deg P}$; the printed $m^2$ is correct only for $\deg P=2$ [?].
\end{coursenote}

Remarks. (1) The massless dispersion relations are obtained likewise with $p\in N^+$ and $m=0$. (2) $P$ enters the decomposition twice: first in deciding what, for a given observer, is the split into purely temporal and purely spatial directions, second in the dispersion relation itself; to play both roles $P$ must be bi-hyperbolic and energy-distinguishing. (3) Once the relation is solved for $E$ as a series in the spatial momentum,
\[
  E=\sum_{i=0}^\infty c^{\alpha_1\dots\alpha_i}p_{\alpha_1}\cdots p_{\alpha_i},\qquad c=m,
\]
one can hardly tell, from the typically infinite list of coefficients $c,c^\alpha,c^{\alpha\beta},\dots$, whether $P$ was bi-hyperbolic and energy-distinguishing, i.e.\ whether the relation belongs to causal matter with a definite energy sign. And bounding the coefficients from experiment makes little sense without the observer frames that give $E$ and $p_\alpha\epsilon^\alpha$ their meaning. The paper therefore has serious doubts about such phenomenological attempts.

\section{Lecture VII: Superluminality}\label{sec:c29-L7}

\obg{Lecture VII} Whenever $\deg P>2$, massive particles can move superluminally. Causality is not endangered, since the definitions rule that out, but one asks why such particles are not routinely seen. The answer given: kinematics permits a superluminal massive particle to shed massless particles and so slow down until it is infraluminal.

\begin{theorem}[Vacuum Cerenkov-type radiation]\label{thm:c29-cerenkov}
\obg{Lecture VII, eq.~(1)}
A positive energy massive particle of momentum $p$ cannot, for kinematic reasons, emit a positive energy massless particle at $x$ exactly when $p$ lies in the cone
\[
  L_x^{-1}\bigl(C^\#_x\bigr)\subseteq C
\]
in cotangent space.
\end{theorem}

Remarks. (1) The cone $L^{-1}_x(C^\#_x)$ holds exactly the momenta of \emph{infraluminal} massive particles, because the slowest light travels along the boundary of $C^\#$; every other positive energy massive particle outruns the slowest light, and in some geometries even the fastest. (2) For a Lorentzian metric ($\deg P=2$) one has $L^{-1}_x(C^\#_x)=C_x$: positive energy massive particles are slower than light and cannot radiate a massless particle in vacuo. The familiar statement that vacuum Cerenkov radiation does not occur therefore holds only for Lorentzian spacetimes. (3) The proof is in [A] (R\"atzel, Rivera, Schuller) and is not reproduced in the lectures.

\paragraph{Reduction to 1+1 dimensions.} To see what lies in store for superluminal particles without cumbersome algebra the paper considers $1+1$ dimensions; the extension to $3+1$ has no conceptual difficulty. The paper's figures show, in a cotangent plane, the cones of massless momenta together with the mass shell of mass $m$ and the construction of distinguished points on it; I describe them.

Because the Gauss map sends the inner massless cones of the cotangent space to outer cones in the tangent space and vice versa, the \emph{inner} cone corresponds to \emph{fast} massless particles and the \emph{outer} one to \emph{slow} massless particles.

Construction. Let $A_1,A_2$ be the two points where the \emph{slow} massless momentum cone touches the mass shell. A massive momentum of higher energy can emit a massless particle moving at the slow speed, because the slow cone through that point lies below the shell. Therefore the shell points between $A_1$ and $A_2$ are exactly the stable ones, and the rays from the origin through $A_1$ and $A_2$ delimit the cone $L^{-1}(C^\#)$ of the theorem. Next let $C$ be the intersection of the slow massless cone centred at $A_2$ with the mass shell: a momentum between $A_1$ and $C$ can emit at most one massless particle, since the remaining massive momentum then falls between $A_2$ and $A_1$, where it cannot decay further (a similar point exists on the left half). Finally let $B$ (and a similar point on the left) be the intersection of the slow massless cone centred at the lowest-energy point of the shell with the shell. It marks where the outgoing massive particle stops running in the same direction as the ingoing one: for ingoing momenta between $B$ and $C$ the outgoing one points the other way.

\begin{coursenote}[Names of the points]
The letter $C$ is used by the paper both for the hyperbolicity cone and for this point of the diagram; the energies $E_{1a},E_{1b},E_2$ below are those of the points $A_1$, $B$, $C$ in this order of increasing energy (I read it so from the text, as the figure is not copied) [?].
\end{coursenote}

Resulting pattern of arrival at a detector, for massive particles emitted by a source with momentum $p$:
\begin{itemize}
  \item $m<E<E_{1a}$ (from the rest mass up to the first energy at which emission of a massless particle is possible): the particle is stable and arrives unchanged;
  \item $E_{1a}<E<E_{1b}$ (the particles are necessarily superluminal; above $E_{1b}$ the massive particle reverses direction after an emission): one detects either the superluminal particle with unchanged momentum (no massless particle was radiated) or an infraluminal massive particle whose momentum follows from the diagram of the decay as a function of the emitted momentum;
  \item $E_{1b}<E<E_2$ ($E_2$ is the largest energy for which at most one emission is possible): only unchanged superluminal particles can arrive, because an emitting particle would have its momentum reversed and end up in the stable region $m<E<E_{1a}$.
\end{itemize}
So the 1+1 model predicts a band of energies in which \emph{no} superluminal particle reaches the detector. The quantum decay rate on tensorial spacetimes is computable but outside the scope of the lectures.

\section{Lecture VIII: Spacetime dynamics}\label{sec:c29-L8}

\obg{Lecture VIII} The goal is a law that moves geometric initial data from one $P$-spacelike hypersurface to the next, so that sweeping over $M$ rebuilds a bi-hyperbolic, energy-distinguishing dispersion relation everywhere. Hypersurfaces with hyperbolic conormals are described by embedding maps; one computes how functionals of the embedding change under normal and tangential deformations, and writes these changes with linear deformation operators. Canonical dynamics for the $(\deg P,0)$-tensor field $P$ then means representing the commutator algebra of these operators on the phase space of spatial geometries. In this book the Lorentzian case is chapters~\ref{ch:c18} and~\ref{ch:c19}; the story in words is \cref{sec:c26-dynamics}.

\begin{definition}[New concepts: Normal and tangent vectors of an embedded hypersurface]\label{def:c29-embed}
Let $X:\Sigma\hookrightarrow M$ be a smooth embedding of a manifold $\Sigma$ of dimension $\dim M-1$ (coordinates $y^\alpha$) into $M$ (coordinates $x^a$). The $\dim M$ vectors
\[
  T(y):=L_{X(y)}(n(y)),\qquad e_1(y):=\frac{\partial X^a(y)}{\partial y^1}\frac\partial{\partial x^a},\ \dots,\ e_{\dim M-1}(y):=\frac{\partial X^a(y)}{\partial y^{\dim M-1}}\frac\partial{\partial x^a}
\]
form a basis of each $T_{X(y)}M$ along $\Sigma$, where the conormal $n(y)$ is uniquely fixed by
\[
  n\in C,\qquad P(n)=1,\qquad n(e_\alpha)=0\quad\text{for all }\alpha=1,\dots,\dim M-1.
\]
$T(y)$ is the (spacetime) \emph{normal vector}, $e_1,\dots,e_{\dim M-1}$ are the (spacetime) \emph{tangent vectors} to $X(\Sigma)$. The uniquely determined dual basis is $n(y),\epsilon^1(y),\dots,\epsilon^{\dim M-1}(y)$.
\end{definition}

Remarks. (1) The geometry shows up only in the normal vector $L(n)$, in three ways: through the cone $C$ (which decides whether $\Sigma$ is an admissible initial surface), through the normalisation of the conormal, otherwise free up to positive rescaling, and through the Legendre map taking conormal to normal vector. (2) The embedding determines the \emph{purely spatial geometry}
\[
  P^{\alpha_1\dots\alpha_I}(y)[X]:=P\bigl(\underbrace{\epsilon^{\alpha_1}(y),\dots,\epsilon^{\alpha_I}(y)}_{I},\underbrace{n(y),\dots,n(y)}_{\deg P-I}\bigr),\qquad I=2,\dots,\deg P,
\]
on the hypersurface: the spatial geometry \emph{detected by point particles}. (3) There is also a spatial geometry \emph{seen by fields}: all the non-zero hypersurface tensors obtained by feeding $\epsilon^\alpha,n$ and $e_\alpha,T$ into the covector and vector slots of $G$. How many there are depends on the valence and symmetries of $G$. For an area metric ($G^{abcd}=G^{cdab}=-G^{bacd}$) in $d$ dimensions one gets three fields:
\begin{align*}
  G^{\alpha_1\alpha_2}(y)[X]&=G\bigl(n(y),\epsilon^{\alpha_1}(y),n(y),\epsilon^{\alpha_2}(y)\bigr),\\
  G^{\alpha_1\alpha_2\alpha_3}(y)[X]&=G\bigl(n(y),\epsilon^{\alpha_1}(y),\epsilon^{\alpha_2}(y),\epsilon^{\alpha_3}(y)\bigr),\\
  G^{\alpha_1\alpha_2\alpha_3\alpha_4}(y)[X]&=G\bigl(\epsilon^{\alpha_1}(y),\epsilon^{\alpha_2}(y),\epsilon^{\alpha_3}(y),\epsilon^{\alpha_4}(y)\bigr).
\end{align*}
The paper focuses on the point-particle geometry; the principle for fields is the same. (4) The quantities in (2) and (3) are tensors on $\Sigma$ and \emph{functionals} of the embedding $X$. Similar quantities for $I=0$ and $I=1$ are constantly $1$ and $0$, because of the definition of $n$. (5) Given a spacetime $(M,G,C^\#)$ with its $P$ and an embedding $X$, one can study how functionals of $X$, e.g.\ the geometric data of (2) and (3), change under deformations of the hypersurface. Normal and tangential deformations are encoded by a hypersurface scalar $N$ and a hypersurface vector $N^\alpha e_\alpha$; the change of a functional $F$ of $X$ is, to linear order, $\mathcal H(N)F$ and $\mathcal D(N^\alpha e_\alpha)F$, with the \emph{normal} and \emph{tangential deformation operators}
\[
  \mathcal H(N)=\int_\Sigma d^{d-1}y\,N(y)\,T^a(y)\frac\delta{\delta X^a(y)},\qquad\mathcal D(N^\alpha e_\alpha)=\int_\Sigma d^{d-1}y\,N^\alpha(y)\,e^a_\alpha(y)\frac\delta{\delta X^a(y)}.
\]

\begin{theorem}[Hypersurface deformation algebra]\label{thm:c29-algebra}
On functionals defined on an initial data hypersurface the deformation operators obey
\begin{align*}
  [\mathcal H(N),\mathcal H(M)]&=-\mathcal D\bigl((\deg P-1)\,P^{\alpha\beta}(M\partial_\beta N-N\partial_\beta M)\partial_\alpha\bigr),\\
  [\mathcal D(N^\alpha\partial_\alpha),\mathcal H(M)]&=-\mathcal H(N^\alpha\partial_\alpha M),\\
  [\mathcal D(N^\alpha\partial_\alpha),\mathcal D(M^\beta\partial_\beta)]&=-\mathcal D\bigl((N^\beta\partial_\beta M^\alpha-M^\beta\partial_\beta N^\alpha)\partial_\alpha\bigr).
\end{align*}
\end{theorem}
\begin{recall}[Deformation algebra in the lecture]
\cref{thm:c26-algebra}: there only the structure was known, here is the full formula with the factor $\deg P-1$ and the inverse spatial tensor $P^{\alpha\beta}$. For $\deg P=2$ and $P^{\alpha\beta}=g^{\alpha\beta}$ it is the algebra of canonical gravity, chapters~\ref{ch:c18}, \ref{ch:c19}.
\end{recall}
The paper gives no proof of the theorem. In the paper $M$ denotes both the manifold and a function in $\mathcal H(M)$.

Remarks. (1) The background geometry enters the algebra only through the normal vectors $T(y)$, obtained from the canonical conormals by the Legendre map, which is determined by $P$. Only the first relation depends on the spacetime geometry; the other two are independent of it. (2) These operators say how functionals change on moving to a nearby hypersurface, \emph{provided the whole spacetime is known}. Dynamics for the spatial geometry must instead predict the spatial geometry on the nearby surface from that on the first, with no spacetime given. That can work only if missing knowledge of the surroundings is made up by extra data on the surface, the \emph{canonical momenta}
\[
  \hat\pi_{\alpha_1\alpha_2},\ \hat\pi_{\alpha_1\alpha_2\alpha_3},\ \dots,\ \hat\pi_{\alpha_1\alpha_2\dots\alpha_{\deg P}},
\]
associated with the purely spatial geometric data
\[
  \hat P^{\alpha_1\alpha_2},\ \hat P^{\alpha_1\alpha_2\alpha_3},\ \dots,\ \hat P^{\alpha_1\alpha_2\dots\alpha_{\deg P}},
\]
now no longer induced from a known spacetime but independent tensors on $\Sigma$ (marked with a hat to avoid confusion). The space of the fields $(P^A,\pi_A)$ is the \emph{geometric phase space}; requiring that $\pi_A$ be canonically conjugate to $P^A$ fixes the Poisson bracket
\[
  \{\hat F,\hat G\}=\int_\Sigma dy\,\Bigl[\frac{\delta\hat F}{\delta\hat P^A(y)}\frac{\delta\hat G}{\delta\hat\pi_A(y)}-\frac{\delta\hat G}{\delta\hat P^A(y)}\frac{\delta\hat F}{\delta\hat\pi_A(y)}\Bigr],
\]
with $\hat F,\hat G$ functionals of the phase space variables; a capital index $A$ stands for the whole collection of index groups $(\alpha_1\alpha_2,\alpha_1\alpha_2\alpha_3,\dots,\alpha_1\dots\alpha_{\deg P})$, and a sum over $A$ is a sum over all of these. (3) For the geometry seen by fields one builds the phase space analogously. For an area metric in $d$ dimensions the initial data are the three tensor fields $\hat G^{\alpha_1\alpha_2},\hat G^{\alpha_1\alpha_2\alpha_3},\hat G^{\alpha_1\alpha_2\alpha_3\alpha_4}$ with conjugate momenta $\hat\gamma_{\alpha_1\alpha_2},\hat\gamma_{\alpha_1\alpha_2\alpha_3},\hat\gamma_{\alpha_1\alpha_2\alpha_3\alpha_4}$, and the Poisson brackets are constructed accordingly.

\begin{definition}[New concepts: Dynamics: Hamiltonian and representation of the algebra]\label{def:c29-dynamics}
Time evolution of observables on the geometric phase space is generated by the Hamiltonian
\[
  H=\int_\Sigma d^{d-1}y\,\bigl[\hat{\mathcal H}(N)+\hat{\mathcal D}(N^\alpha e_\alpha)\bigr],
\]
where the \emph{superhamiltonian} $\hat{\mathcal H}$ and the \emph{supermomentum} $\hat{\mathcal D}$ are functionals of $(P^A,\pi_A)$ whose Poisson brackets reproduce the deformation algebra:
\begin{align*}
  \{\hat{\mathcal H}(A),\hat{\mathcal H}(B)\}&=\hat{\mathcal D}\bigl((\deg P-1)\,\hat P^{\alpha\beta}(A\partial_\beta B-B\partial_\beta A)\partial_\alpha\bigr),\\
  \{\hat{\mathcal D}(A^\alpha\partial_\alpha),\hat{\mathcal H}(B)\}&=\hat{\mathcal H}(A^\alpha\partial_\alpha B),\\
  \{\hat{\mathcal D}(A^\alpha\partial_\alpha),\hat{\mathcal D}(B^\beta\partial_\beta)\}&=\hat{\mathcal D}\bigl((A^\beta\partial_\beta B^\alpha-B^\beta\partial_\beta A^\alpha)\partial_\alpha\bigr).
\end{align*}
\end{definition}

Remarks. (1) This is what dynamics means: geometric data on an initial surface are evolved so that, put on a nearby surface (obtained by the deformation fields $N$, $N^\alpha e_\alpha$), and then again on a nearby one, and so on, all the spatial geometries together form a whole spacetime geometry. The representation requirement ensures that the data on the single surfaces are really the spatial geometries induced by that spacetime, with all required properties. (2) The third relation fixes the \emph{supermomentum} of the point-particle geometry uniquely:
\[
  \hat{\mathcal D}(\vec N)=\sum_{I=2}^{\deg P}\int_\Sigma dy\,N^\beta(y)\Bigl[\partial_\beta\hat P^{\alpha_1\dots\alpha_I}\hat\pi_{\alpha_1\dots\alpha_I}+I\,\partial_{\alpha_1}\bigl(\hat P^{\alpha_1\dots\alpha_I}\hat\pi_{\alpha_2\dots\alpha_I\beta}\bigr)\Bigr].
\]
(3) From the second relation the localised superhamiltonian $\hat{\mathcal H}(y):=\hat{\mathcal H}(\delta_y)$ is a scalar density of weight one and splits into a local and a non-local part,
\[
  \hat{\mathcal H}(y)[\hat P,\hat\pi]=\hat{\mathcal H}_{\rm local}(y)[\hat P](\hat\pi)+\hat{\mathcal H}_{\rm non\text{-}local}(y)[\hat P,\hat\pi],
\]
whose non-local part is directly determined:
\[
  \hat{\mathcal H}_{\rm non\text{-}local}(y)=\sum_{I=2}^{\deg P}\Bigl[(\deg P-I)\,\partial_\beta\bigl(\hat P^{\beta\alpha_1\dots\alpha_I}\hat\pi_{\alpha_1\dots\alpha_I}\bigr)-(\deg P-1)\,I\,\partial_\beta\bigl(\hat P^{\alpha_2\dots\alpha_I}\hat P^{\alpha_1\beta}\hat\pi_{\alpha_1\dots\alpha_I}\bigr)\Bigr](y).
\]
Only the local part remains to be found from the first relation.

(4) With the supermomentum and the non-local part known, the first relation is a condition \emph{quadratic} in the local superhamiltonian. Remarkably it is equivalent to a \emph{linear} condition on the Legendre transform of $\hat{\mathcal H}_{\rm local}$ in the variables $\hat\pi_A$,
\[
  L(y)[\hat P](K):=\hat\pi_A(y)[\hat P](K)\,K^A(y)-\hat{\mathcal H}(y)_{\rm local}[\hat P]\bigl(\hat\pi[\hat P](K)\bigr),\qquad K^A(y):=\frac{\partial\hat{\mathcal H}(y)_{\rm local}}{\partial\pi_A(y)}.
\]
With these definitions the first relation becomes the linear functional differential equation
\begin{multline*}
  0=-\frac{\delta L(x)}{\delta\hat P^A(y)}K^A(y)+\partial_{y^\zeta}\Bigl[\frac{\delta L(x)}{\delta\hat P^A(y)}M^{A\zeta}(y)\Bigr]-\frac{\partial L(x)}{\partial K^A(x)}K^B(x)\,Q_B{}^{A\beta}(x)\,\partial_\beta\delta_x(y)\\
  +\frac{\partial L(x)}{\partial K^A(x)}\Bigl[R^{A\mu\nu}(x)\,\partial^2_{\mu\nu}\delta_x(y)-S^{A\mu}(x)\,\partial_\mu\delta_x(y)\Bigr]-(x\leftrightarrow y),
\end{multline*}
whose coefficients $R^{A\mu\nu}$, $Q_A{}^{B\mu}$, $M^{A\beta}$ contain only the configuration variables, not their derivatives,
\begin{align*}
  R^{\alpha_1\dots\alpha_I\mu\nu}&=I(\deg P-1)\,P^{(\beta|(\alpha_1}P^{\alpha_2\dots\alpha_I)|\mu)},\\
  M^{\alpha_1\dots\alpha_I\beta}&=-(\deg P-I)\hat P^{\beta\alpha_1\dots\alpha_I}+I(\deg P-1)\hat P^{\beta(\alpha_1}\hat P^{\alpha_2\dots\alpha_I)},\\
  Q_{\alpha_1\dots\alpha_K}{}^{\beta_1\dots\beta_I\mu}&=\delta^K_{I+1}(\deg P-I)\,\delta^{\mu\beta_1\dots\beta_I}_{(\alpha_1\dots\alpha_{I+1})}-\delta^K_2\,I(\deg P-1)\hat P^{(\beta_2\dots\beta_I}\delta^{\beta_1)\mu}_{(\alpha_1\alpha_2)}\\
  &\quad-\delta^K_{I-1}\,I(\deg P-1)\hat P^{\mu(\beta_1}\delta^{\beta_2\dots\beta_I)}_{\alpha_1\dots\alpha_{I-1}},
\end{align*}
while $S^{A\mu}$ contains also their first partial derivatives:
\begin{align*}
  S^{\alpha_1\dots\alpha_I\mu}&=I(\deg P-1)(\deg P-2)\,P^{\mu\gamma(\alpha_1}\partial_\gamma P^{\alpha_2\dots\alpha_I)}+I(\deg P-1)(\deg P-I)\,\partial_\gamma P^{\gamma(\alpha_1}P^{\alpha_2\dots\alpha_I)\mu}\\
  &\quad+(\deg P-1)\,P^{\mu\gamma}\partial_\gamma P^{\alpha_1\dots\alpha_I}-I(I-1)(\deg P-1)^2\,\partial_\gamma P^{\gamma(\alpha_1}P^{\alpha_2\dots\alpha_{I-1}}P^{\alpha_I)\mu}.
\end{align*}
(5) The linearity of this equation in $L$ permits the power series ansatz
\[
  L(x)[\hat P](K)=\sum_{i=0}^\infty G(x)\bigl(\hat P,\partial\hat P,\partial\partial\hat P,\dots\bigr)_{A_1\dots A_i}K^{A_1}(x)\cdots K^{A_i}(x),
\]
and the search for the local superhamiltonian, hence for the gravitational dynamics, becomes the search for the coefficients $G_{A_1\dots A_N}$ obeying the equations that this ansatz produces. They are given in the theorem below.

\begin{coursenote}[Printed formulas of the coefficients]
The coefficients $R$, $M$, $Q$, $S$ and the equations (I)--(VIII) below are reproduced as printed, symbol by symbol. They contain what look like misprints in the preprint (for example, in $R^{\alpha_1\dots\alpha_I\mu\nu}$ the free index appears as $\beta$ on the right where $\nu$ is meant; in the $\{\mathcal H,\mathcal H\}$ bracket the capital letters $A,B$ are smearing functions and not the index $A$ of the phase space; the factor ``$NN!$'' in (III) is presumably $N\,N!$). They have not been rederived here, and should be checked against the final paper by Giesel, Schuller, Witte and Wohlfarth before being used in a calculation [?]. The equations are shown as formulas to be read, not to be solved by hand.
\end{coursenote}

\begin{theorem}[Linear PDE system for the gravitational Lagrangian]\label{thm:c29-pde}
The coefficients $G_A,G_{AB},G_{ABC},\dots$ of the local superhamiltonian of the point-particle geometry (which together with the known non-local part and the supermomentum give the full gravitational dynamics)
\begin{enumerate}
  \item[(a)] involve nothing beyond the variables $\hat P^A$ and their partial derivatives of order at most two;
  \item[(b)] are fixed entirely by the equations expressing the representation of the deformation algebra, namely
\end{enumerate}
\begin{align*}
\text{(I)}\quad 0={}&2\partial_\mu(G_AR^{A\beta\mu})+2G_AS^{A\beta}-2\partial_\mu\Bigl(\frac{\partial G}{\partial\partial_{(\mu|}\hat P^A}M^{A|\beta)}\Bigr)-4\partial_\mu\Bigl(\frac{\partial G}{\partial\partial^2_{(\mu|\nu}\hat P^A}\partial_\nu M^{A|\beta)}\Bigr)\\
&+2M^{A\beta}\frac{\partial G}{\partial\hat P^A}+2\partial_\mu M^{A\beta}\frac{\partial G}{\partial\partial_\mu\hat P^A}+2\partial^2_{\mu\nu}M^{A\beta}\frac{\partial G}{\partial\partial^2_{\mu\nu}\hat P^A},\\
\text{(II)}\quad 0={}&(N+1)!\,G_{AB_1\dots B_N}R^{A\alpha\beta}-N!\,\frac{\partial G_{B_1\dots B_N}}{\partial\partial_{(\beta|}\hat P^A}M^{A|\alpha)}-2N!\,\frac{\partial G_{B_1\dots B_N}}{\partial\partial^2_{(\beta|\gamma}\hat P^A}\partial_\gamma M^{A|\alpha)}\\
&-(N-2)(N-1)!\,\frac{\partial G_{B_1\dots B_{N-1}}}{\partial\partial^2_{\alpha\beta}\hat P^{B_N}},\\
\text{(III)}\quad 0={}&(N+1)!\,G_{AB_1\dots B_N}S^{A\alpha}+(N-1)!\sum_{a=1}^N\frac{\partial G_{B_1\dots\widetilde{B_a}\dots B_N}}{\partial\partial_\alpha\hat P^{B_a}}-2(N-1)!\,\partial_\gamma\frac{\partial G_{B_1\dots B_{N-1}}}{\partial\partial^2_{\alpha\gamma}\hat P^{B_N}}\\
&+N!\,\frac{G_{B_1\dots B_N}}{\partial\hat P^A}M^{A\alpha}+N!\,\frac{\partial G_{B_1\dots B_N}}{\partial\partial_\gamma\hat P^A}\partial_\gamma M^{A\alpha}+N!\,\frac{\partial G_{B_1\dots B_N}}{\partial\partial^2_{\gamma\delta}\hat P^A}\partial^2_{\gamma\delta}M^{A\alpha}\\
&-N\,N!\,Q_{(B_1}{}^{M\alpha}G_{B_2\dots B_N)M},
\end{align*}
which also include the two \emph{symmetry conditions}
\begin{align*}
\text{(IV)}\quad&\frac{\partial G_{B_1\dots\widetilde{B_a}\dots B_N}}{\partial\partial^2_{\gamma_1\gamma_2}\hat P^{B_a}}=\frac{\partial G_{B_1\dots B_{N-1}}}{\partial\partial^2_{\gamma_1\gamma_2}\hat P^{B_N}}\qquad\text{for all }N\ge1,\ a=1,\dots,N,\\
\text{(V)}\quad&0=\frac{\partial G_{B_1\dots B_N}}{\partial\partial^2_{(\alpha\beta|}\hat P^A}M^{A|\gamma)}\qquad\text{for all }N\ge0,
\end{align*}
and the three \emph{invariance conditions}
\begin{align*}
\text{(VI)}\quad0={}&\sum_{I=2}^{\deg P}I\,\hat P^{\alpha_2\dots\alpha_I(\sigma}\frac{\partial G_{B_1\dots B_N}}{\partial\partial^2_{\mu\nu)}\hat P^{\alpha_2\dots\alpha_I\rho}},\\
\text{(VII)}\quad0={}&\sum_{I=2}^{\deg P}\Bigl[I\,\hat P^{\alpha_2\dots\alpha_I(\mu}\frac{\partial G_{B_1\dots B_N}}{\partial\partial_{\nu)}\hat P^{\alpha_2\dots\alpha_I\rho}}-\partial_\rho\hat P^{\alpha_1\dots\alpha_I}\frac{\partial G_{B_1\dots B_N}}{\partial\partial^2_{\mu\nu}\hat P^{\alpha_1\dots\alpha_I}}\\
&\qquad\quad+2I\,\partial_\sigma\hat P^{\alpha_2\dots\alpha_I(\mu}\frac{\partial G_{B_1\dots B_N}}{\partial\partial^2_{\nu)\sigma}\hat P^{\alpha_2\dots\alpha_I\rho}}\Bigr],\\
\text{(VIII)}\quad&\sum_{I=2}^{\deg P}\Bigl[I\,\hat P^{\rho\beta_2\dots\beta_I}\frac{\partial G_{B_1\dots B_N}}{\partial\hat P^{\beta_2\dots\beta_I\mu}}+I\,\partial_\gamma\hat P^{\rho\beta_2\dots\beta_I}\frac{\partial G_{B_1\dots B_N}}{\partial\partial_\gamma\hat P^{\beta_2\dots\beta_I\mu}}-\partial_\mu\hat P^{\beta_1\dots\beta_I}\frac{\partial G_{B_1\dots B_N}}{\partial\partial_\rho\hat P^{\beta_1\dots\beta_I}}\\
&\qquad+I\,\partial_{\gamma\delta}\hat P^{\rho\beta_2\dots\beta_I}\frac{\partial G_{B_1\dots B_N}}{\partial\partial_{\gamma\delta}\hat P^{\beta_2\dots\beta_I\mu}}-2\,\partial_{\mu\gamma}\hat P^{\beta_1\dots\beta_I}\frac{\partial G_{B_1\dots B_N}}{\partial\partial_{\rho\gamma}\hat P^{\beta_1\dots\beta_I}}\Bigr]\\
&=-\delta^\rho_\mu G_{B_1\dots B_N}-n_1\,\delta^\rho_{(\beta^{(1)}_1}G_{\beta^{(1)}_2\dots\beta^{(1)}_{n_1})\mu B_2\dots B_N}-\dots-n_N\,G_{B_1\dots B_{N-1}\mu(\beta^{(N)}_2\dots\beta^{(N)}_{n_N}}\delta^\rho_{\beta^{(N)}_1)},
\end{align*}
where $n_i$ is the number of small indices contained in the capital index $B_i=\beta^{(i)}_1\dots\beta^{(i)}_{n_i}$.
\end{theorem}
\begin{recall}[The system in the lecture]
\cref{thm:c26-pde}: ``six equations and five sequences of equations''. The paper has three equations (I)--(III), with two symmetry conditions (IV), (V) and three invariance conditions (VI)--(VIII); the lecture counted them differently, and I did not try to match the two counts.
\end{recall}
The paper gives the result without derivation; the derivation is promised in [B].

Remarks. (1) To find a gravity theory for the point-particle geometry is to solve these linear PDEs and then express $\hat\pi_A$ through the Legendre-dual variables $K^A$: a purely \emph{mathematical problem}. The field geometry is treated the same way. (2) For $\deg P=2$, $\dim M=4$ the system is comparatively easy to solve and the only non-zero coefficients are
\begin{align*}
  G&=(2\kappa)^{-1}\sqrt g\,(R-2\lambda),\\
  G_{\alpha\beta}&=\rho\sqrt g\,\bigl(R_{\alpha\beta}-\tfrac12g_{\alpha\beta}R\bigr)+\sigma\sqrt g\,g_{\alpha\beta},\\
  G_{\alpha\beta\mu\nu}&=(16\kappa)^{-1}\sqrt g\,\bigl[g_{\alpha\mu}g_{\beta\nu}+g_{\beta\mu}g_{\alpha\nu}-2g_{\alpha\beta}g_{\mu\nu}\bigr],
\end{align*}
where $R_{\alpha\beta}$ and $R$ are the Ricci tensor and scalar of $P^{\alpha\beta}=g^{\alpha\beta}$, and $\kappa,\lambda,\rho,\sigma$ are four free real integration constants. The last two drop out of the equations of motion and may be set to zero; $\kappa$ and $\lambda$ are the gravitational and cosmological constants, which only experiment can fix. The result is Einstein's general relativity. (3) Open questions are therefore: (i) existence of modified gravity, i.e.\ solvability of the system for even $\deg P>2$; (ii) uniqueness for each such degree; (iii) the explicit solutions, once they exist.
\begin{recall}[Maxwell gives Einstein]
\cref{ex:c26-einstein}: in the lecture the result for Maxwell matter was quoted as $\frac1G\sqrt{-g}(R+\Lambda)$ with undetermined constants; here are the coefficient functions of the expansion with the constants $\kappa,\lambda$ (and $\rho,\sigma$ that drop out).
\end{recall}

\begin{coursenote}[Draft text left in the arXiv source, not in the 44-page paper]
After the end of the document the LaTeX source still contains an unused draft title page and a closing paragraph for the last lecture; they are not in the compiled paper, so they are not part of the chapter above. Content, in my words: the dynamics of $P$ is not a matter of postulating an action with deep intuition; the causal structure of a bi-hyperbolic, energy-distinguishing $P$ already fixes the deformation algebra of the initial data hypersurfaces, and the dynamics must represent it on the phase space of the induced spatial geometries. For Lorentzian geometry this is the argument of the geometrodynamics school around Wheeler (the draft cites two works by keys that are not resolved), which leads to the Einstein--Hilbert action with a free cosmological constant. Hence finding dynamics reduces to the mathematical problem solved by \cref{thm:c29-pde}. The draft also announces that the deformation algebra is only sketched.
\end{coursenote}

\paragraph{Background and further reading in the paper.} The ten specialist texts named: Isham, \emph{Differential Geometry for Physicists}; Stewart, \emph{Advanced General Relativity}; H\"ormander, \emph{Linear partial differential operators}; Hehl and Obukhov, \emph{Foundations of classical electrodynamics}; Beem, Ehrlich and Easley, \emph{Global Lorentzian geometry}; Shen, \emph{Lectures on Finsler geometry}; Perlick, \emph{Ray optics, Fermat's principle and applications to general relativity}; Hassett, \emph{Introduction to algebraic geometry}; Sundermeyer, \emph{Constrained dynamics}; Rockafellar, \emph{Convex analysis}. The paper acknowledges Klaus Mecke for the invitation to the summer school in Obergurgl.

\section*{Key concepts}
\begin{keyconcepts}
\textbf{Retrieval questions.}
\begin{enumerate}
  \item What is a probing matter field, and why is it required to be linear? Which condition on the geometry does well-posedness of the initial value problem give (\cref{def:c29-probing}, \cref{thm:c29-hyp})?
  \item Define a hyperbolic polynomial and its hyperbolicity cone; give the Hurwitz construction of the cone and the cone of $g^{-1}(k,k)$.
  \item What is a locally wavelike solution, and how does one read the dispersion relation $P(x,k)=0$ from it? When is a field equation massive?
  \item Define the dual polynomial. Why is $[DP^\#]$ the partial inverse of $[DP]$, and which condition on $DDP$ is needed (\cref{thm:c29-gauss})?
  \item Write the action of a free massless particle in terms of $P$, and in terms of $P^\#$.
  \item State the three conditions that make a geometry a spacetime with respect to matter dynamics. Which role has each (predictivity, time orientation, agreement on the sign of energy)?
  \item Why must an energy-distinguishing polynomial have even degree (\cref{rem:c29-evenness})?
  \item What is the barrier function, why is the Legendre map invertible, and what does it map? Derive the action $m\int P^*(\dot x)^{1/\deg P}d\tau$.
  \item Why do superluminal massive particles not appear at a detector in a range of energies in $1+1$ dimensions (\cref{thm:c29-cerenkov})?
  \item What are the normal and the tangential deformation operators, their algebra, and which part of it depends on the geometry? What does it mean to represent it on the geometric phase space, and to what does it reduce for $\deg P=2$?
\end{enumerate}
\textbf{Concepts.}
\begin{itemize}
  \item Geometry, probing matter, principal polynomial, hyperbolicity, hyperbolicity cone (\cref{def:c29-geometry}, \cref{def:c29-probing}, \cref{thm:c29-hyp}, \cref{def:c29-cone}).
  \item Locally wavelike solutions, dispersion relation, massless momenta (\cref{def:c29-wavelike}, \cref{def:c29-null}).
  \item Dual polynomial, Gauss and dual Gauss map, massless particle action (\cref{def:c29-dual}, \cref{thm:c29-gauss}, \cref{thm:c29-massless}).
  \item Bi-hyperbolic, time orientation, positive energy cone, energy-distinguishing, spacetime (\cref{def:c29-bihyp}, \cref{def:c29-timeor}, \cref{def:c29-poscone}, \cref{def:c29-endist}, \cref{def:c29-spacetime}).
  \item Mass, barrier function, Legendre map, massive particle action, observer, dispersion relation (\cref{def:c29-massive}, \cref{thm:c29-legendre}, \cref{thm:c29-massivepart}, \cref{def:c29-observer}, \cref{thm:c29-dispersion}).
  \item Superluminality and the absence of vacuum Cerenkov radiation only for $\deg P=2$ (\cref{thm:c29-cerenkov}).
  \item Hypersurface deformation algebra, geometric phase space, superhamiltonian and supermomentum, linear PDE system; Einstein gravity for $\deg P=2$ (\cref{thm:c29-algebra}, \cref{def:c29-dynamics}, \cref{thm:c29-pde}).
\end{itemize}
\end{keyconcepts}

\flushside
\end{paracol}
